\documentclass[11pt]{article}

% ── Geometry and encoding ──────────────────────────────────────────
\usepackage[margin=1in]{geometry}
\usepackage[utf8]{inputenc}
\usepackage[T1]{fontenc}

% ── Mathematics ────────────────────────────────────────────────────
\usepackage{amsmath,amssymb,amsthm}
\usepackage{mathpartir}
\usepackage{stmaryrd}

% ── Code listings ──────────────────────────────────────────────────
\usepackage{listings}
\usepackage{xcolor}

% ── Figures, tables, links ─────────────────────────────────────────
\usepackage{graphicx}
\usepackage{booktabs}
\usepackage{multirow}
\usepackage{hyperref}
\usepackage{cleveref}
\usepackage{enumitem}
\usepackage{float}
\usepackage{caption}
\usepackage{subcaption}

% ── Bibliography ───────────────────────────────────────────────────
\usepackage[numbers,sort&compress]{natbib}

% ── Theorem environments ──────────────────────────────────────────
\newtheorem{theorem}{Theorem}[section]
\newtheorem{lemma}[theorem]{Lemma}
\newtheorem{proposition}[theorem]{Proposition}
\newtheorem{corollary}[theorem]{Corollary}
\newtheorem{definition}[theorem]{Definition}
\newtheorem{example}[theorem]{Example}
\newtheorem{remark}[theorem]{Remark}

% ── JAPL listing style ────────────────────────────────────────────
\definecolor{japlkw}{RGB}{0,100,170}
\definecolor{japlstr}{RGB}{180,60,30}
\definecolor{japlcmt}{RGB}{100,120,100}
\definecolor{japltype}{RGB}{140,50,140}
\definecolor{japlbg}{RGB}{248,248,248}

\lstdefinelanguage{japl}{
  keywords={fn,let,match,with,if,then,else,type,module,import,use,own,ref,
            resource,trait,impl,test,property,forall,loop,while,do,continue,
            opaque,signature,foreign,unsafe,deriving,packed,bench,assert},
  keywordstyle=\color{japlkw}\bfseries,
  keywords=[2]{Int,Float,Float32,Bool,Char,String,Bytes,Unit,Never,
               List,Option,Some,None,Result,Ok,Err,Map,Set,
               Owned,Io,Async,Net,State,Process,Fail,Pure,
               File,Socket,Buffer,Pid,Reply,Config,Path,
               True,False,Permanent,Transient,Temporary,
               ReadError,MutableArray,DbConn},
  keywordstyle=[2]\color{japltype},
  comment=[l]{--},
  commentstyle=\color{japlcmt}\itshape,
  string=[b]",
  stringstyle=\color{japlstr},
  sensitive=true,
  morestring=[b]',
  literate={->}{$\to$}2 {=>}{$\Rightarrow$}2 {|>}{$\triangleright$}2
           {>>}{$\gg$}2 {++}{$\mathbin{+\mkern-5mu+}$}2
           {>=}{$\geq$}2 {<=}{$\leq$}2 {!=}{$\neq$}2,
}

\lstset{
  language=japl,
  basicstyle=\small\ttfamily,
  backgroundcolor=\color{japlbg},
  frame=single,
  rulecolor=\color{black!20},
  numbers=left,
  numberstyle=\tiny\color{gray},
  numbersep=6pt,
  tabsize=2,
  breaklines=true,
  breakatwhitespace=true,
  showstringspaces=false,
  columns=flexible,
  captionpos=b,
  aboveskip=10pt,
  belowskip=10pt,
}

% ── Macros ─────────────────────────────────────────────────────────
\newcommand{\JAPL}{\textsc{Japl}}
\newcommand{\own}{\mathsf{own}}
\newcommand{\rref}{\mathsf{ref}}
\newcommand{\Owned}[1]{\mathsf{Owned}\langle #1 \rangle}
\newcommand{\Io}{\mathsf{Io}}
\newcommand{\Pure}{\mathsf{Pure}}
\newcommand{\Resource}{\mathsf{Resource}}
\newcommand{\Value}{\mathsf{Value}}
\newcommand{\sem}[1]{\llbracket #1 \rrbracket}
\newcommand{\lin}{\mathsf{lin}}
\newcommand{\un}{\mathsf{un}}
\newcommand{\lolli}{\multimap}
\newcommand{\tensor}{\otimes}
\newcommand{\parr}{\mathbin{\bindnasrepma}}
\newcommand{\bang}{{\mathop{!}}}
\newcommand{\why}{{\mathop{?}}}
\newcommand{\GC}{\textsc{gc}}

% ── Title ──────────────────────────────────────────────────────────
\title{%
  \textbf{Mutation Is Local and Explicit:}\\[4pt]
  \large Controlled Effects and Linear Ownership\\
  in the \JAPL{} Programming Language%
}

\author{
  Matthew Long \\
  \textit{The JAPL Research Collaboration} \\
  \textit{YonedaAI Research Collective} \\
  Chicago, IL \\
  \texttt{matthew@yonedaai.com}
}

\date{March 2026}

% ══════════════════════════════════════════════════════════════════
\begin{document}
\maketitle

% ── Abstract ──────────────────────────────────────────────────────
\begin{abstract}
Shared mutable state is the single most prolific source of defects in concurrent software.
We present \JAPL{}'s approach to this problem: a \emph{dual-layer} type system that
separates \emph{pure values}---immutable, garbage-collected, and freely
shared---from \emph{linear resources}---mutable, ownership-tracked, and
deterministically released.  The pure layer handles the overwhelming
majority of application logic with no annotation burden, while the
resource layer gives systems-level control over files, sockets,
buffers, and foreign memory through an ownership discipline rooted in
linear type theory and Girard's linear logic.  We formalize the
dual-layer system, prove type safety (progress and preservation), and
show that well-typed \JAPL{} programs satisfy \emph{resource safety}:
every acquired resource is released exactly once.  A detailed
comparison with Rust's borrow checker, Haskell's \texttt{IO} monad,
Clean's uniqueness types, and several other approaches demonstrates
that \JAPL{}'s design achieves a favorable trade-off between
annotation burden, safety guarantees, and runtime performance.
\end{abstract}

\medskip
\noindent\textbf{Keywords:}
linear types, ownership, mutation, resource safety, effect systems,
substructural type systems, dual-layer runtime

\newpage
\tableofcontents
\newpage

% ══════════════════════════════════════════════════════════════════
\section{Introduction}\label{sec:intro}

The history of programming language design can be read, in large part,
as a long negotiation with mutable state.  Assembly gives the
programmer direct access to every register and memory cell.  C
organizes that access into functions and structs but leaves aliasing
entirely to the programmer's discipline.  Java introduces garbage
collection and access modifiers but permits unrestricted sharing of
mutable objects across threads.  Haskell confines all side effects to
the \texttt{IO} monad, achieving referential transparency at the cost
of pervasive monadic plumbing.  Rust enforces ownership and borrowing
rules that prevent data races at compile time, but its borrow checker
pervades \emph{every} layer of the program, even pure data
transformations that could never cause a data race.

\JAPL{} takes a different position: \textbf{mutation is local and
explicit}.  The language is organized around two sharply delineated
layers:

\begin{enumerate}[label=(\roman*)]
  \item A \textbf{pure layer} where values are immutable, garbage-collected,
        and freely shared.  Most application code---business logic,
        data transformations, protocol handling---lives here.
  \item A \textbf{resource layer} where mutable resources (file handles,
        network sockets, GPU buffers, FFI pointers) are tracked by a
        linear ownership discipline.  Resources must be consumed exactly
        once: used and then released.
\end{enumerate}

This separation is the second of \JAPL{}'s seven core design
principles\footnote{The seven principles are: (1)~Values are primary,
(2)~Mutation is local and explicit, (3)~Concurrency is process-based,
(4)~Failures are normal and typed, (5)~Distribution is a native
concern, (6)~Functions are the unit of composition,
(7)~Runtime simplicity matters as much as type power.} and is the
subject of this paper.

\subsection{The Mutation Problem}

Shared mutable state creates three interrelated difficulties:

\paragraph{Concurrency bugs.}
When two threads can both read and write the same memory location, the
result depends on scheduling---a data race.  Languages without
mutation restrictions must rely on locks, atomic operations, or
transactional memory, each introducing its own class of bugs
(deadlocks, priority inversion, live-lock).

\paragraph{Aliasing and invalidation.}
When multiple references point to the same mutable object, a mutation
through one reference can invalidate assumptions made through another.
The classic example is iterator invalidation: removing an element from
a collection while iterating over it.  More subtle is the
\emph{borrowed reference} problem in systems code: a pointer into a
buffer becomes dangling after the buffer is reallocated.

\paragraph{Reasoning difficulty.}
Equational reasoning---substituting equals for equals---fails in the
presence of mutation.  If \texttt{f(x)} mutates \texttt{x}, then
\texttt{let y = f(x); g(y, x)} is not equivalent to
\texttt{g(f(x), x)}, because the second occurrence of \texttt{x}
observes different state.  This makes refactoring hazardous and
formal verification prohibitively expensive.

\JAPL{}'s dual-layer design addresses all three problems.  Pure values
cannot be mutated, so they are immune to data races, aliasing bugs,
and reasoning failures.  Resources \emph{can} be mutated, but
ownership tracking ensures that at most one reference exists at any
time, restoring the ability to reason locally about mutation effects.

\subsection{Contributions}

This paper makes the following contributions:

\begin{enumerate}
  \item We present a formal account of \JAPL{}'s dual-layer type system,
        connecting it to Girard's linear logic and the theory of
        substructural type systems (\Cref{sec:formal,sec:dual-layer}).
  \item We define the ownership semantics for the resource layer,
        including consumption, borrowing, region-based inference, and
        uniqueness inference (\Cref{sec:ownership}).
  \item We formalize \JAPL{}'s resource types and prove that well-typed
        programs satisfy \emph{resource safety}: every acquired
        resource is eventually released exactly once (\Cref{sec:resource-types,sec:proofs}).
  \item We describe the interaction between the two layers, including
        the borrowing protocol and escape hatches (\Cref{sec:interaction}).
  \item We provide a detailed comparison with six other language
        designs (\Cref{sec:comparison}) and describe an implementation
        strategy for the dual-layer runtime (\Cref{sec:implementation}).
\end{enumerate}

\subsection{Paper Organization}

\Cref{sec:background} surveys related work.  \Cref{sec:formal}
develops the formal framework: linear logic, substructural type
systems, and categorical semantics.  \Cref{sec:dual-layer} presents
\JAPL{}'s dual-layer model.  \Cref{sec:ownership} details ownership
semantics.  \Cref{sec:resource-types} defines resource types.
\Cref{sec:interaction} describes layer interaction.
\Cref{sec:comparison} compares with related languages.
\Cref{sec:implementation} discusses implementation.
\Cref{sec:proofs} contains formal proofs.
\Cref{sec:discussion} reflects on practical trade-offs.
\Cref{sec:conclusion} concludes.


% ══════════════════════════════════════════════════════════════════
\section{Background and Related Work}\label{sec:background}

\subsection{Substructural Type Systems}

The structural rules of a type system determine how hypotheses in a
typing context may be used.  Classical logic (and the simply-typed
lambda calculus) admits three structural rules: \emph{weakening}
(a hypothesis may be ignored), \emph{contraction} (a hypothesis may
be duplicated), and \emph{exchange} (hypotheses may be reordered).
Dropping one or more of these rules yields a \emph{substructural}
type system~\cite{walker2005substructural}:

\begin{definition}[Substructural type systems]\label{def:substructural}
Let $\Gamma$ denote a typing context.  The four principal substructural
disciplines are:
\begin{enumerate}[label=(\alph*)]
  \item \textbf{Linear}: No weakening, no contraction.  Every variable
        must be used exactly once.
  \item \textbf{Affine}: No contraction, but weakening is allowed.
        Every variable may be used at most once.
  \item \textbf{Relevant}: No weakening, but contraction is allowed.
        Every variable must be used at least once.
  \item \textbf{Ordered}: No weakening, no contraction, no exchange.
        Variables must be used exactly once, in order of introduction.
\end{enumerate}
\end{definition}

\JAPL{}'s resource layer employs \emph{linear} types: each resource
must be consumed exactly once.  The pure layer uses standard
structural types (with all three rules), since immutable values can
be freely shared, duplicated, or ignored.

\subsection{Linear Logic}

Girard's linear logic~\cite{girard1987linear} provides the
proof-theoretic foundation for linear type systems.  The key
connectives are:

\begin{itemize}
  \item $A \lolli B$ (linear implication): consuming one $A$ produces
        one $B$.
  \item $A \tensor B$ (tensor / multiplicative conjunction): both $A$
        and $B$ are available simultaneously, each used exactly once.
  \item $\bang A$ (of course / exponential): an unlimited supply of
        $A$, recoverable as many times as needed.
  \item $A \,\&\, B$ (with / additive conjunction): a choice between
        $A$ and $B$, where only one is consumed.
  \item $A \oplus B$ (plus / additive disjunction): exactly one of $A$
        or $B$ is provided.
\end{itemize}

The exponential $\bang A$ is critical: it bridges linear and
unrestricted worlds.  In \JAPL{}, the pure layer corresponds to the
$\bang$-fragment of linear logic: every pure value implicitly carries
an exponential modality, meaning it can be freely duplicated and
discarded.  The resource layer corresponds to the linear fragment
proper.

\subsection{Rust's Ownership and Borrowing}\label{sec:bg-rust}

Rust~\cite{rust2021} pioneered the use of ownership and borrowing as a
practical mechanism for memory safety without garbage collection.
Rust's key rules are:

\begin{enumerate}
  \item Each value has a single owner.
  \item Ownership can be transferred (moved).
  \item References can be borrowed: either one mutable reference
        \emph{or} any number of immutable references, but not both
        simultaneously.
  \item Values are dropped (destructors run) when the owner goes out of
        scope.
\end{enumerate}

The RustBelt project~\cite{jung2017rustbelt} formalized Rust's type
system using the Iris separation logic framework, proving that
well-typed Rust programs (even those using \texttt{unsafe}) satisfy
memory safety when the unsafe blocks are locally sound.

\JAPL{} differs from Rust in a fundamental architectural decision:
Rust applies ownership to \emph{all} values, requiring lifetime
annotations even for purely functional data transformations.  \JAPL{}
restricts ownership to the resource layer, leaving the pure layer
free of ownership concerns.

\subsection{Clean's Uniqueness Types}\label{sec:bg-clean}

Clean~\cite{barendsen1996uniqueness,plasmeijer2001clean} uses
\emph{uniqueness types} to ensure that a value has at most one
reference.  If a value is unique, it can be destructively updated
in place without violating referential transparency, because no
other reference can observe the mutation.

Clean's uniqueness attribute \texttt{*} is propagated by the type
checker.  For example, \texttt{*File} denotes a unique file handle
that can be read or written.  Uniqueness types are closely related to
linear types: a unique value must be used exactly once (or
threaded through a computation that returns a new unique value).

\JAPL{}'s resource layer is inspired by Clean's approach but extends
it with explicit \texttt{own} and \texttt{ref} qualifiers, region
inference, and integration with an algebraic effect system.

\subsection{Linear Haskell}\label{sec:bg-linear-haskell}

Linear Haskell~\cite{bernardy2018linear} extends Haskell with linear
arrow types $a \lolli b$, allowing functions to specify that an
argument must be consumed exactly once.  The key insight is
\emph{retrofitting}: linear types are added to an existing language
without breaking backward compatibility, by making linearity a
property of function arrows rather than types.

\JAPL{} takes a complementary approach.  Rather than adding linearity
to arrows, \JAPL{} adds linearity to a distinguished class of
\emph{types} (resource types).  This makes the boundary between linear
and unrestricted code syntactically visible.

\subsection{Region-Based Memory Management}\label{sec:bg-regions}

Tofte and Talpin~\cite{tofte1997region} introduced region-based memory
management, where objects are allocated into lexically scoped regions
and deallocated en masse when the region goes out of scope.  The
Cyclone language~\cite{jim2002cyclone,grossman2002region} brought
regions to a C-like setting with static safety guarantees.

Region inference automatically determines which region each allocation
belongs to, minimizing annotation burden.  \JAPL{} uses region-like
reasoning in its ownership inference pass: the compiler determines the
lifetime of each resource and inserts release operations at the
appropriate scope boundaries.

\subsection{Capability-Based Security}\label{sec:bg-capabilities}

The capability model of security~\cite{dennis1966programming,
miller2006robust} organizes access control around unforgeable tokens
(capabilities) rather than identity-based access control lists.  A
process can only perform an operation if it holds the appropriate
capability.

\JAPL{} applies this principle to resource access: a file can only be
read if the function holds an \texttt{Owned<File>} or \texttt{ref File}
capability.  The effect system (\texttt{Io}, \texttt{Net}) further
restricts which capabilities can be exercised in a given context.

\subsection{Austral, Roc, and Koka}\label{sec:bg-austral-roc-koka}

Several recent languages explore design points closely related to
\JAPL{}'s.

\paragraph{Austral.}
Austral~\cite{austral2023} is a systems language with linear types in
which every type is either \emph{free} (unrestricted, may be copied
and discarded) or \emph{linear} (must be consumed exactly once).  The
dichotomy mirrors \JAPL{}'s pure/resource split, but Austral applies
it within a single layer: there is no garbage-collected pure heap, and
all memory management is explicit.  \JAPL{}'s dual-layer design adds a
GC-managed pure layer for ergonomics, and borrowing via $\rref\;T$ for
non-consuming reads, which Austral lacks.

\paragraph{Roc.}
Roc is a functional language that targets application-level programming
and uses a capability-based effect system in which all side effects
flow through \emph{platforms}.  Like \JAPL{}, Roc keeps the vast
majority of code pure and pushes effects to the boundary.  Unlike
\JAPL{}, Roc does not expose linear or ownership types to the
programmer; resource management is delegated entirely to the platform
layer.

\paragraph{Koka.}
Koka~\cite{leijen2017koka} pioneered row-typed algebraic effects with
automatic effect inference.  \JAPL{}'s effect system draws on Koka's
design, but Koka does not provide linear types or an ownership
discipline for resources.  The combination of row-typed effects
\emph{and} linear resource tracking is a distinctive feature of
\JAPL{}'s design.

\subsection{Algebraic Effects and Handlers}

Algebraic effects~\cite{plotkin2009handlers,bauer2015programming} provide a
structured alternative to monads for describing side effects.  An
effect signature declares a set of operations; an effect handler
provides implementations for those operations.

\JAPL{}'s effect system (\texttt{Io}, \texttt{State[s]},
\texttt{Fail[e]}, etc.) draws on algebraic effects.  Effects form a
commutative monoid under composition, and effect handlers serve as
natural transformations from effectful computations to pure
results~\cite{kammar2013handlers}.  The interaction between effects and
linear types is a key contribution of this paper: resource operations
are effectful (\texttt{with Io}), and the linear discipline ensures
that effect handlers cannot accidentally duplicate or discard
resources.


% ══════════════════════════════════════════════════════════════════
\section{Formal Framework}\label{sec:formal}

This section develops the type-theoretic and logical foundations of
\JAPL{}'s dual-layer system.

\subsection{Linear Type Theory}\label{sec:linear-type-theory}

We begin with a standard presentation of the linear lambda
calculus~\cite{wadler1990linear,wadler1993taste}.

\begin{definition}[Linear types]\label{def:linear-types}
The types of the linear lambda calculus are given by the grammar:
\[
  A, B ::= \alpha \mid A \lolli B \mid A \tensor B \mid \bang A
           \mid A \,\&\, B \mid A \oplus B \mid \mathbf{1} \mid \mathbf{0}
\]
where $\alpha$ ranges over type variables, $\lolli$ is linear
function space, $\tensor$ is multiplicative conjunction, $\bang$ is
the exponential, $\&$ is additive conjunction, $\oplus$ is additive
disjunction, $\mathbf{1}$ is the unit of $\tensor$, and $\mathbf{0}$
is the empty type.
\end{definition}

\begin{definition}[Linear typing context]\label{def:linear-context}
A \emph{linear context} $\Delta$ is a multiset of typing assumptions
$x : A$ where each variable $x$ appears at most once.  An
\emph{unrestricted context} $\Gamma$ contains assumptions $x : A$
where $x$ may be used arbitrarily many times.  A \emph{mixed context}
is a pair $\Gamma; \Delta$.
\end{definition}

The typing judgment $\Gamma; \Delta \vdash e : A$ asserts that
expression $e$ has type $A$ under unrestricted context $\Gamma$ and
linear context $\Delta$, where \emph{every} assumption in $\Delta$ is
consumed exactly once.

\begin{definition}[Core typing rules]\label{def:core-rules}
The following rules govern the linear lambda calculus with exponentials:

\begin{mathpar}
  \inferrule*[right=Var-Lin]
    { }
    {\Gamma; x : A \vdash x : A}

  \inferrule*[right=Var-Un]
    {x : A \in \Gamma}
    {\Gamma; \cdot \vdash x : A}

  \inferrule*[right=$\lolli$-I]
    {\Gamma; \Delta, x : A \vdash e : B}
    {\Gamma; \Delta \vdash \lambda x.\, e : A \lolli B}

  \inferrule*[right=$\lolli$-E]
    {\Gamma; \Delta_1 \vdash e_1 : A \lolli B \\
     \Gamma; \Delta_2 \vdash e_2 : A}
    {\Gamma; \Delta_1, \Delta_2 \vdash e_1\; e_2 : B}

  \inferrule*[right=$\tensor$-I]
    {\Gamma; \Delta_1 \vdash e_1 : A \\
     \Gamma; \Delta_2 \vdash e_2 : B}
    {\Gamma; \Delta_1, \Delta_2 \vdash (e_1, e_2) : A \tensor B}

  \inferrule*[right=$\tensor$-E]
    {\Gamma; \Delta_1 \vdash e : A \tensor B \\
     \Gamma; \Delta_2, x : A, y : B \vdash e' : C}
    {\Gamma; \Delta_1, \Delta_2 \vdash
      \mathbf{let}\; (x, y) = e\; \mathbf{in}\; e' : C}

  \inferrule*[right=$\bang$-I]
    {\Gamma; \cdot \vdash e : A}
    {\Gamma; \cdot \vdash \bang e : \bang A}

  \inferrule*[right=$\bang$-E]
    {\Gamma; \Delta_1 \vdash e_1 : \bang A \\
     \Gamma, x : A; \Delta_2 \vdash e_2 : B}
    {\Gamma; \Delta_1, \Delta_2 \vdash
      \mathbf{let}\; \bang x = e_1\; \mathbf{in}\; e_2 : B}
\end{mathpar}
\end{definition}

The $\bang$-I rule requires that the expression $e$ use \emph{no}
linear resources: only unrestricted values can be promoted to the
exponential.  This is the formal counterpart of \JAPL{}'s rule that
pure values cannot capture resources.

\subsection{Connection to Girard's Linear Logic}\label{sec:girard}

The Curry-Howard correspondence extends to linear logic: proofs
correspond to programs, propositions to types, and cut-elimination to
evaluation~\cite{girard1987linear,wadler2012propositions}.

\begin{proposition}[Curry-Howard for linear logic]\label{prop:curry-howard}
Under the linear Curry-Howard correspondence:
\begin{enumerate}[label=(\roman*)]
  \item $A \lolli B$ corresponds to a function consuming one $A$ to produce one $B$.
  \item $A \tensor B$ corresponds to a pair where both components must be used.
  \item $\bang A$ corresponds to a value that may be duplicated or discarded.
  \item $A \,\&\, B$ corresponds to a lazy pair (choose one projection).
  \item $A \oplus B$ corresponds to a tagged union (one alternative is provided).
\end{enumerate}
\end{proposition}

In \JAPL{}'s dual-layer system, the Curry-Howard correspondence
manifests directly: pure-layer types correspond to the
$\bang$-fragment (all types implicitly carry $\bang$), while
resource-layer types correspond to the linear fragment (no implicit
$\bang$).

\subsection{Categorical Semantics}\label{sec:categorical}

The categorical semantics of linear types provides a rigorous
foundation for the dual-layer design.

\begin{definition}[Symmetric monoidal closed category]\label{def:smcc}
A \emph{symmetric monoidal closed category} (SMCC)
$(\mathcal{C}, \tensor, I, \lolli)$ consists of a category $\mathcal{C}$
equipped with a bifunctor $\tensor : \mathcal{C} \times \mathcal{C}
\to \mathcal{C}$, a unit object $I$, and an internal hom
$\lolli$ such that $- \tensor A \dashv A \lolli -$ for every object
$A$.
\end{definition}

\begin{definition}[Linear-non-linear model]\label{def:lnl}
A \emph{linear-non-linear (LNL) model}~\cite{benton1995mixed} consists
of:
\begin{enumerate}[label=(\roman*)]
  \item A Cartesian closed category $\mathcal{C}$ (the ``non-linear'' or ``intuitionistic'' part).
  \item A symmetric monoidal closed category $\mathcal{L}$ (the ``linear'' part).
  \item A symmetric monoidal adjunction $F \dashv G : \mathcal{L}
        \to \mathcal{C}$ where $F : \mathcal{C} \to \mathcal{L}$
        is a symmetric monoidal functor.
\end{enumerate}
The exponential $\bang$ is recovered as the comonad $FG$ on $\mathcal{L}$.
\end{definition}

\begin{theorem}[Benton's LNL correspondence~\cite{benton1995mixed}]\label{thm:lnl}
The category of LNL models is equivalent to the category of models
of intuitionistic linear logic (ILL) with an exponential $\bang$
satisfying the standard structural rules (weakening, contraction,
dereliction, promotion).
\end{theorem}

This result is directly relevant to \JAPL{}:

\begin{remark}[\JAPL{} as an LNL model]\label{rem:japl-lnl}
\JAPL{}'s dual-layer system instantiates the LNL framework:
\begin{itemize}
  \item $\mathcal{C}$ is the category of pure \JAPL{} types, which is
        Cartesian closed (it has product types, function types, and
        a terminal object \texttt{Unit}).
  \item $\mathcal{L}$ is the category of resource \JAPL{} types, which
        is symmetric monoidal closed (tensor is pair-of-resources,
        linear function is ownership-transferring function).
  \item The functor $F : \mathcal{C} \to \mathcal{L}$ embeds pure
        values into the resource layer (a pure value can always be
        used where a resource is expected, since it carries an
        implicit $\bang$).
  \item The functor $G : \mathcal{L} \to \mathcal{C}$ ``forgets''
        linearity---this corresponds to the \texttt{freeze} or
        \texttt{snapshot} operations that convert a resource to an
        immutable value.
\end{itemize}
\end{remark}

\subsection{Compact Closed Categories and Duality}\label{sec:compact-closed}

In the full linear logic setting, the multiplicative fragment admits
a notion of \emph{duality}: every type $A$ has a dual $A^\perp$,
and the category of types forms a $\ast$-autonomous
category~\cite{barr1991autonomous}.

\begin{definition}[$\ast$-Autonomous category]\label{def:star-auto}
A $\ast$-\emph{autonomous category} is a symmetric monoidal closed
category $(\mathcal{C}, \tensor, I, \lolli)$ equipped with a
dualizing object $\bot$ such that the canonical morphism
$A \to (A \lolli \bot) \lolli \bot$ is an isomorphism for every
object $A$.
\end{definition}

While \JAPL{} does not expose full classical linear logic to the
programmer, the duality principle manifests in the resource layer's
\emph{consumption pattern}: acquiring a resource of type $R$ creates
an obligation (dual) to release it.  The linear type system ensures
that this obligation is always discharged.


% ══════════════════════════════════════════════════════════════════
\section{JAPL's Dual-Layer Model}\label{sec:dual-layer}

This section presents the core design of \JAPL{}'s two-layer type and
memory system.

\subsection{Design Philosophy}

The dual-layer model is motivated by an empirical observation:
\textbf{most code does not need mutation}.  Business logic, data
transformations, protocol parsing, serialization, routing---all of
these are naturally expressed as pure functions over immutable values.
Mutation is needed only at the boundaries: reading files, writing
sockets, managing GPU memory, interacting with foreign libraries.

By restricting ownership tracking to the resource layer, \JAPL{}
achieves two goals simultaneously:

\begin{enumerate}
  \item \textbf{Ergonomics}: The vast majority of code is free of
        ownership annotations, lifetime parameters, and borrow-checker
        errors.
  \item \textbf{Safety}: The code that \emph{does} interact with
        mutable resources is subject to a rigorous linear discipline
        that prevents resource leaks, use-after-free, and
        double-free.
\end{enumerate}

\subsection{The Pure Layer}\label{sec:pure-layer}

The pure layer encompasses all immutable values: algebraic data types,
records, tuples, lists, maps, strings, closures, and numeric types.

\begin{definition}[Pure types]\label{def:pure-types}
The pure types of \JAPL{} are given by:
\[
  \tau ::= \alpha \mid \texttt{Int} \mid \texttt{Float} \mid \texttt{Bool}
         \mid \texttt{String} \mid \texttt{Unit}
         \mid \tau_1 \times \tau_2 \mid \tau_1 + \tau_2
         \mid \tau_1 \to \tau_2 \mid \texttt{List}[\tau]
         \mid \{ \ell_1 : \tau_1, \ldots, \ell_n : \tau_n \mid \rho \}
\]
where $\alpha$ ranges over type variables and $\rho$ ranges over row
variables.
\end{definition}

Pure values have the following properties:

\begin{itemize}
  \item \textbf{Immutable}: Once created, a pure value never changes.
  \item \textbf{GC-managed}: Memory is reclaimed by a generational,
        per-process garbage collector.
  \item \textbf{Freely shareable}: Pure values can be passed to any
        number of functions, stored in any number of data structures,
        and sent to any number of processes.
  \item \textbf{Structurally typed}: Two values with the same structure
        have the same type (row polymorphism for records).
\end{itemize}

\begin{lstlisting}[caption={Pure-layer code: no ownership concerns},label={lst:pure-layer}]
-- Pure values: immutable, GC-managed, freely shareable
let data = [1, 2, 3, 4, 5]
let copy = data  -- sharing is fine; data is immutable

-- Pure function: transforms data, no side effects
fn transform(data: List[Int]) -> List[Int] =
  List.map(data, fn x -> x * 2)

-- Composition of pure functions
fn process_order(order: Order) -> Invoice =
  order
  |> validate
  |> calculate_totals
  |> apply_discounts
  |> generate_invoice
\end{lstlisting}

\subsection{The Resource Layer}\label{sec:resource-layer}

The resource layer encompasses all mutable, owned values: file handles,
network sockets, database connections, GPU buffers, FFI pointers, and
mutable arrays.

\begin{definition}[Resource types]\label{def:resource-types}
A resource type is introduced by the \texttt{resource} keyword:
\[
  R ::= \texttt{resource}\; T
\]
Resource types are \emph{linear}: each value of a resource type must
be consumed exactly once in every execution path.
\end{definition}

\begin{definition}[Ownership qualifiers]\label{def:ownership-qualifiers}
\JAPL{} provides two ownership qualifiers:
\begin{enumerate}[label=(\alph*)]
  \item $\own\; T$: Exclusive ownership.  The holder can read, write,
        and release the resource.  Corresponds to the linear type $T$.
  \item $\rref\; T$: Borrowed reference.  The holder can read the
        resource but cannot write, release, or transfer it.
        Corresponds to $\bang T$ restricted to a scope.
\end{enumerate}
\end{definition}

\begin{lstlisting}[caption={Resource-layer code: explicit ownership},label={lst:resource-layer}]
resource File
resource Socket

-- Opening a file acquires ownership (manual management via let)
fn read_config(path: Path) -> Result[Config, ReadError] with Io =
  let file = File.open(path, Read)?
  let contents = File.read_all(file)?
  File.close(file)  -- ownership consumed; compile error if omitted
  parse_config(contents)

-- Ownership transfer: send consumes the socket, returns it
fn send(socket: Owned[Socket], bytes: Bytes)
    -> (Owned[Socket], Int) with Io =
  let n = Socket.write(socket, bytes)
  (socket, n)
\end{lstlisting}

\subsection{The Split in Practice}

The boundary between pure and resource layers is clean and
syntactically visible.  Consider a web server:

\begin{lstlisting}[caption={The dual-layer split in a web server},label={lst:web-server}]
-- Pure layer: all business logic
fn route(req: Request) -> Response =
  match req.method, req.path with
  | Get, "/users/" ++ id -> get_user(id)
  | Post, "/users"       -> create_user(req.body)
  | _                    -> not_found()

fn get_user(id: String) -> Response =
  match UserRepo.find(id) with
  | Some(user) -> Response.json(200, User.to_json(user))
  | None       -> Response.json(404, error_json("not found"))

-- Resource layer: IO boundary
fn serve(listener: Owned[TcpListener])
    -> Never with Io, Process[ServerMsg] =
  let (listener, conn) = Tcp.accept(listener)
  let _ = Process.spawn(fn -> handle_conn(conn))
  serve(listener)

fn handle_conn(conn: Owned[TcpSocket]) -> Unit with Io =
  let (conn, req_bytes) = Tcp.read(conn)
  let req = parse_request(req_bytes)  -- pure
  let resp = route(req)               -- pure
  let conn = Tcp.write(conn, serialize_response(resp))
  Tcp.close(conn)
\end{lstlisting}

In this example, the pure functions \texttt{route}, \texttt{get\_user},
\texttt{parse\_request}, and \texttt{serialize\_response} make up the
majority of the code and require no ownership annotations.  Only
\texttt{serve} and \texttt{handle\_conn}, which interact with TCP
sockets, operate in the resource layer.

\subsection{Formal Typing Judgment}

The dual-layer system uses a combined typing judgment:

\begin{definition}[Dual-layer typing judgment]\label{def:dual-judgment}
The judgment
\[
  \Gamma; \Delta \vdash_\varepsilon e : \tau
\]
asserts that expression $e$ has type $\tau$ under unrestricted
(pure) context $\Gamma$, linear (resource) context $\Delta$, and
with effect $\varepsilon \in \{\Pure, \Io, \ldots\}$.

The rules ensure:
\begin{enumerate}[label=(\roman*)]
  \item Variables in $\Gamma$ may be used any number of times.
  \item Variables in $\Delta$ must each be used exactly once.
  \item Resource-producing operations (\texttt{File.open}, etc.) add
        bindings to $\Delta$.
  \item Resource-consuming operations (\texttt{File.close}, etc.)
        remove bindings from $\Delta$.
  \item At every join point (if-then-else, match), the remaining
        $\Delta$ must be identical in all branches.
\end{enumerate}
\end{definition}


% ══════════════════════════════════════════════════════════════════
\section{Ownership Semantics}\label{sec:ownership}

\subsection{Owned$\langle T \rangle$}

The type $\Owned{T}$ represents exclusive ownership of a resource of
type $T$.  The owner has three rights:

\begin{enumerate}
  \item \textbf{Use}: Perform operations on the resource (read, write).
  \item \textbf{Transfer}: Pass ownership to another function or process.
  \item \textbf{Release}: Deallocate the resource, releasing its
        underlying system resource.
\end{enumerate}

These rights are mutually exclusive over time: once ownership is
transferred or the resource is released, the original holder loses all
rights.

\begin{definition}[Ownership transfer]\label{def:transfer}
A function that accepts an $\Owned{T}$ parameter \emph{consumes} the
caller's ownership.  If the function returns an $\Owned{T}$, ownership
is transferred back to the caller.  The typing rule is:
\begin{mathpar}
  \inferrule*[right=Own-Transfer]
    {\Gamma; \Delta, x : \Owned{T} \vdash_\varepsilon e : \tau}
    {\Gamma; \Delta \vdash_\varepsilon
      \lambda(x : \Owned{T}).\, e : \Owned{T} \lolli_\varepsilon \tau}
\end{mathpar}
\end{definition}

\subsection{Consumption and Return Patterns}\label{sec:consumption}

\JAPL{} supports several patterns for resource consumption:

\paragraph{Terminal consumption.}
The resource is released and not returned:
\begin{lstlisting}[caption={Terminal consumption: File.close},label={lst:terminal}]
fn close_file(file: Owned[File]) -> Unit with Io =
  File.close(file)
  -- file is consumed; cannot be used after this point
\end{lstlisting}

\paragraph{Threaded consumption.}
The resource is used and then returned with updated state:
\begin{lstlisting}[caption={Threaded consumption: read then return},label={lst:threaded}]
fn read_line(file: Owned[File])
    -> (Owned[File], String) with Io =
  let (file, line) = File.read_line(file)
  (file, line)
\end{lstlisting}

\paragraph{Fork-join consumption.}
Ownership is temporarily split (borrowing) and then rejoined:
\begin{lstlisting}[caption={Fork-join via borrowing},label={lst:fork-join}]
fn copy_file(src: Owned[File], dst: Owned[File])
    -> (Owned[File], Owned[File]) with Io =
  let bytes = File.read_all(src)
  let dst = File.write_all(dst, bytes)
  (src, dst)
\end{lstlisting}

\subsection{Borrowing}\label{sec:borrowing}

Borrowing allows temporary read-only access to a resource without
transferring ownership.

\begin{definition}[Borrow]\label{def:borrow}
A \emph{borrow} of type $\rref\; T$ grants read-only access to a
resource of type $T$.  The borrow is valid only within the lexical
scope of the borrowing expression.  The typing rule is:
\begin{mathpar}
  \inferrule*[right=Borrow]
    {\Gamma; \Delta, x : \Owned{T} \vdash_\varepsilon e : \tau \\
     x \notin \mathrm{FV}(e) \text{ as owned}}
    {\Gamma; \Delta, x : \Owned{T} \vdash_\varepsilon
      \mathbf{borrow}\; x\; \mathbf{as}\; y\; \mathbf{in}\; [y/x]e : \tau}
\end{mathpar}
where $y : \rref\; T$ is the borrowed reference, and $x$ remains
available as $\Owned{T}$ after the borrow scope exits.
\end{definition}

\begin{remark}[Borrow desugaring]\label{rem:borrow-desugar}
The surface syntax $\mathbf{borrow}\; x\; \mathbf{as}\; y\;
\mathbf{in}\; e$ desugars to the core calculus form
$\mathsf{borrow}(\mathsf{res}(a),\; \lambda y.\, e)$, where $x$ is
bound to $\mathsf{res}(a)$ and the body $e$ is wrapped in a lambda
that receives the borrowed value $y$.  The core reduction rule
(Definition~\ref{def:reduction}, rule \textsc{Borrow}) then reduces
$\mathsf{borrow}(\mathsf{res}(a), f)$ to $(f\; v,\;
\mathsf{res}(a))$, applying the closure to the live value $v$ while
preserving the resource handle.
\end{remark}

\begin{lstlisting}[caption={Borrowing: peek at a buffer without consuming it},label={lst:borrow}]
fn peek(buf: ref Buffer) -> Byte =
  Buffer.get(buf, 0)

fn process(buf: Owned[Buffer]) -> (Owned[Buffer], Byte) with Io =
  let byte = peek(ref buf)  -- borrow for read-only access
  (buf, byte)               -- buf is still owned here
\end{lstlisting}

The key constraint is that borrows cannot \emph{outlive} the owner:
the borrowed reference $\rref\; T$ is valid only within the scope
where it is created.  This is enforced statically by the compiler.

\subsection{Region-Based Inference}\label{sec:region-inference}

\JAPL{} uses region-based inference~\cite{tofte1997region} to
automatically determine resource lifetimes and insert release
operations.

\begin{definition}[Region]\label{def:region}
A \emph{region} $r$ is an abstract lifetime associated with a lexical
scope.  Every resource allocation is assigned a region, and the
resource is released when the region ends.  Regions form a partial
order: if scope $r_1$ is nested inside scope $r_2$, then
$r_1 \sqsubseteq r_2$ (i.e., $r_1$ ends before $r_2$).
\end{definition}

The compiler's region inference pass operates in three phases:

\begin{enumerate}
  \item \textbf{Constraint generation}: For each resource allocation,
        generate a region variable.  For each use, generate a
        constraint that the resource's region must contain the use
        site.  For each ownership transfer, propagate the region.
  \item \textbf{Constraint solving}: Solve the region constraints to
        find the minimal region for each resource.
  \item \textbf{Insertion}: Insert release operations at region
        boundaries where resources go out of scope.
\end{enumerate}

\begin{example}[Automatic release]\label{ex:auto-release}
In the following code, the compiler infers that \texttt{conn} must
live at least as long as the scope of \texttt{process\_query}:
\begin{lstlisting}[caption={Region inference inserts Db.close automatically},label={lst:region}]
fn process_query(url: String, sql: String)
    -> Result[Rows, DbError] with Io =
  use conn = Db.connect(url)?
  let result = Db.query(ref conn, sql)?
  -- compiler inserts: Db.close(conn) here
  Ok(result)
\end{lstlisting}
The \texttt{use} binding signals that \texttt{conn} is a linear
resource; the compiler ensures it is released at the end of its
region.
\end{example}

\subsection{Uniqueness Inference}\label{sec:uniqueness-inference}

Beyond region inference, \JAPL{} performs \emph{uniqueness inference}
to determine when a value has a single reference, enabling in-place
mutation optimizations.

\begin{definition}[Uniqueness]\label{def:uniqueness}
A value is \emph{unique} if it has exactly one reference in the
program state.  A unique immutable value can be safely updated
in-place (destructive update), because no other reference can observe
the change.
\end{definition}

\begin{proposition}[Uniqueness optimization]\label{prop:uniqueness-opt}
If a pure value $v$ is unique at a record-update site
$\{ v \mid \ell = e \}$, the compiler may perform the update in-place,
avoiding allocation.  This is semantically equivalent to
creating a fresh copy with the field changed.
\end{proposition}

\begin{proof}
Since $v$ has a single reference, no other expression can observe the
difference between in-place update and copy-on-write.  By
parametricity, the program cannot distinguish the two implementations.
\end{proof}

This optimization is particularly important for the common pattern
of threading state through recursive functions:

\begin{lstlisting}[caption={Uniqueness inference enables in-place update},label={lst:uniqueness}]
fn server_loop(state: ServerState) -> Never with Process[Msg] =
  let msg = Process.receive()
  -- If state is unique (no other references), the update
  -- { state | count = state.count + 1 } is performed in-place.
  let new_state = { state | count = state.count + 1 }
  server_loop(new_state)
\end{lstlisting}


% ══════════════════════════════════════════════════════════════════
\section{Resource Types}\label{sec:resource-types}

\subsection{Declaring Resource Types}

Resource types are declared with the \texttt{resource} keyword:

\begin{lstlisting}[caption={Resource type declarations},label={lst:resource-decl}]
resource File
resource Socket
resource DbConnection
resource GpuBuffer
resource MutableArray[a]
\end{lstlisting}

The \texttt{resource} keyword instructs the compiler to apply linear
typing rules to values of this type.  A resource cannot be silently
dropped, duplicated, or aliased.

\subsection{Resource Lifecycle}\label{sec:lifecycle}

Every resource follows a lifecycle:

\begin{enumerate}
  \item \textbf{Acquisition}: A resource is created via a constructor
        (e.g., \texttt{File.open}, \texttt{Tcp.connect}).  The
        constructor returns $\Owned{R}$.
  \item \textbf{Use}: The resource is operated on via functions that
        take $\Owned{R}$ or $\rref\; R$.  Functions that take
        $\Owned{R}$ must return it (threaded pattern) or consume it
        (terminal pattern).
  \item \textbf{Release}: The resource is released via a destructor
        (e.g., \texttt{File.close}, \texttt{Tcp.disconnect}).  The
        destructor consumes $\Owned{R}$ and returns no resource.
\end{enumerate}

\begin{definition}[Resource safety]\label{def:resource-safety}
A program is \emph{resource-safe} if, for every resource acquired
during execution, the resource is eventually released exactly once.
\end{definition}

\begin{theorem}[Well-typed programs are resource-safe]\label{thm:resource-safety}
If $\Gamma; \cdot \vdash_\varepsilon e : \tau$ (i.e., the program
type-checks with an empty linear context), then execution of $e$
is resource-safe.
\end{theorem}

We defer the proof to \Cref{sec:proofs}.

\subsection{The \texttt{use} Binding}\label{sec:use-binding}

The \texttt{use} keyword introduces a linear binding with
automatic release semantics:

\begin{lstlisting}[caption={The \texttt{use} binding},label={lst:use-binding}]
fn read_entire_file(path: String) -> Result[String, IoError] with Io =
  use file = File.open(path, Read)?
  let contents = File.read_all(file)?
  -- File.close(file) is automatically inserted by the compiler
  -- at the end of the `use` scope
  Ok(contents)
\end{lstlisting}

The \texttt{use} binding desugars to a try-finally pattern: even if
the body raises an error (via the \texttt{?} operator or a crash),
the resource is released.

\begin{definition}[\texttt{use} desugaring]\label{def:use-desugar}
The binding $\texttt{use}\; x = e_1\; \texttt{in}\; e_2$ desugars to:
\[
  \mathbf{let}\; x = e_1\; \mathbf{in}\;
  \mathbf{try}\; e_2\; \mathbf{finally}\; \mathit{release}(x)
\]
where $\mathit{release}$ is the destructor associated with the resource
type of $x$.
\end{definition}

\subsection{Parametric Resource Types}

Resources can be parametric:

\begin{lstlisting}[caption={Parametric resource types},label={lst:parametric-resource}]
resource MutableArray[a]

fn sort_in_place(arr: own MutableArray[Int])
    -> own MutableArray[Int] =
  let len = MutableArray.length(arr)
  loop i = 0, arr = arr while i < len do
    loop j = i + 1, arr = arr while j < len do
      if MutableArray.get(arr, j) < MutableArray.get(arr, i) then
        let arr = MutableArray.swap(arr, i, j)
        continue(j + 1, arr)
      else
        continue(j + 1, arr)
    continue(i + 1, arr)
  arr
\end{lstlisting}

The linearity constraint applies to the outer resource container, not
to its type parameter.  The elements of a \texttt{MutableArray[Int]}
are pure \texttt{Int} values and can be freely copied.

\subsection{Freezing: Resource to Value Conversion}\label{sec:freezing}

A resource can be \emph{frozen} to produce an immutable value,
transferring data from the resource layer to the pure layer:

\begin{lstlisting}[caption={Freezing a buffer into immutable bytes},label={lst:freeze}]
fn build_message() -> Bytes with Io =
  let buf = Buffer.alloc(1024)          -- resource layer
  let buf = Buffer.write(buf, 0, header)
  let buf = Buffer.write(buf, 64, payload)
  Buffer.freeze(buf)                    -- pure layer: immutable Bytes
  -- buf is consumed; only the frozen Bytes value remains
\end{lstlisting}

\begin{definition}[Freeze]\label{def:freeze}
A \emph{freeze} operation has type $\Owned{R} \lolli \tau$, consuming
the resource and producing a pure value.  After freezing, the resource
is released and the returned value lives in the pure layer (managed
by GC).
\end{definition}

Freezing corresponds to the functor $G : \mathcal{L} \to \mathcal{C}$
in the LNL model (\Cref{def:lnl}): it maps a linear resource to an
unrestricted pure value.


% ══════════════════════════════════════════════════════════════════
\section{Interaction Between Layers}\label{sec:interaction}

\subsection{Pure Functions Using Resources}

Pure functions cannot \emph{own} resources, but they can \emph{borrow}
them.  A function that borrows a resource via $\rref\; T$ does not
participate in the ownership protocol; it simply reads the resource's
current state.

\begin{lstlisting}[caption={Pure function borrowing a resource},label={lst:pure-borrow}]
-- This function is pure except for the borrow
fn buffer_checksum(buf: ref Buffer) -> Int =
  let bytes = Buffer.to_bytes_view(buf)  -- read-only view
  Bytes.fold(bytes, 0, fn acc, b -> acc + b)

fn validate_and_send(buf: Owned[Buffer], sock: Owned[Socket])
    -> (Owned[Buffer], Owned[Socket]) with Io =
  let sum = buffer_checksum(ref buf)  -- borrow for pure computation
  if sum > 0 then
    let (sock, _) = Socket.send(sock, Buffer.to_bytes(ref buf))
    (buf, sock)
  else
    (buf, sock)
\end{lstlisting}

\subsection{Ownership Transfer Between Processes}

When a resource is sent to another process, ownership transfers:

\begin{lstlisting}[caption={Ownership transfer between processes},label={lst:process-transfer}]
fn producer() -> Unit with Io, Process[ProducerMsg] =
  let buf = Buffer.alloc(4096)
  let buf = Buffer.write(buf, 0, generate_data())
  -- Ownership transfers to the consumer process
  Process.send(consumer_pid, ProcessBuffer(buf))
  -- buf is consumed; using it here is a compile error

fn consumer() -> Never with Io, Process[ConsumerMsg] =
  match Process.receive() with
  | ProcessBuffer(buf) ->
      -- We now own buf
      let data = Buffer.freeze(buf)
      process_data(data)
      consumer()
\end{lstlisting}

\subsection{The Escape Hatch}\label{sec:escape-hatch}

For interoperability with C libraries and low-level systems code,
\JAPL{} provides an \texttt{unsafe} escape hatch that allows
bypassing linearity checks:

\begin{lstlisting}[caption={The unsafe escape hatch},label={lst:unsafe}]
-- unsafe block: linearity is not checked
fn raw_mmap(fd: Owned[File], size: Int) -> Owned[MappedRegion] with Io =
  unsafe {
    let ptr = foreign_mmap(File.raw_fd(ref fd), size)
    MappedRegion.from_raw(ptr, size)
  }
  -- Outside unsafe: normal linearity rules apply
\end{lstlisting}

The \texttt{unsafe} keyword has several properties:

\begin{enumerate}
  \item It must appear explicitly in the source code, making it
        searchable and auditable.
  \item It does not propagate: a function containing \texttt{unsafe}
        can still present a safe interface to its callers.
  \item It is tracked by the effect system: \texttt{unsafe} blocks
        require the \texttt{Unsafe} effect, which cannot be silently
        introduced.
\end{enumerate}

\subsection{Effect System Interaction}

Resource operations are effectful: they carry the \texttt{Io} effect.
The effect system and the ownership system cooperate:

\begin{definition}[Effect-ownership interaction]\label{def:effect-ownership}
The following invariants hold:
\begin{enumerate}[label=(\roman*)]
  \item Resource acquisition requires an effect (typically $\Io$).
  \item Resource release requires the same effect as acquisition.
  \item Pure functions ($\varepsilon = \Pure$) cannot acquire or
        release resources; they can only borrow.
  \item Effect handlers cannot duplicate or discard resources: the
        handler's continuation must preserve the linear context.
\end{enumerate}
\end{definition}

\begin{proposition}[Effect handlers preserve linearity]\label{prop:effect-linear}
If an effect handler $h$ transforms computation
$\Gamma; \Delta \vdash_\varepsilon e : \tau$ into
$\Gamma; \Delta' \vdash_{\varepsilon'} e' : \tau'$,
then $\Delta' = \Delta$ (the linear context is unchanged).
\end{proposition}

\begin{proof}
By induction on the structure of the effect handler.  The handler
provides implementations for effect operations, but each operation
implementation must preserve the linear context: it receives a
continuation that expects the same linear resources, and it must
pass exactly those resources to the continuation.  Since the handler
cannot forge or destroy linear resources (this is enforced by the
$\bang$-I rule, which prevents linear resources from appearing in
unrestricted positions), the linear context is preserved.
\end{proof}


% ══════════════════════════════════════════════════════════════════
\section{Comparison with Related Languages}\label{sec:comparison}

\subsection{Overview}

\begin{table}[H]
\centering
\small
\caption{Comparison of mutation control strategies across languages}
\label{tab:comparison}
\begin{tabular}{@{}lllll@{}}
\toprule
\textbf{Language} & \textbf{Mutation model} & \textbf{Resource mgmt} &
  \textbf{Annotation} & \textbf{GC} \\
\midrule
\JAPL{} & Dual-layer & Linear ownership & Low (pure layer) & Hybrid \\
Rust & Ownership+borrow & Ownership+RAII & High (everywhere) & None \\
Go & Unrestricted & GC + \texttt{defer} & None & Full GC \\
Haskell & IO monad & GC + finalizers & Medium (monadic) & Full GC \\
Erlang & Copy everything & Per-process GC & None & Per-proc GC \\
OCaml & Mutable refs & GC & Low & Full GC \\
Clean & Uniqueness types & Uniqueness & Medium & GC \\
\bottomrule
\end{tabular}
\end{table}

\subsection{Rust: Full Borrow Checker}\label{sec:cmp-rust}

Rust applies its ownership and borrowing discipline to \emph{every}
value, including purely functional data.  This provides strong
guarantees---no GC is needed, and memory safety is ensured at compile
time---but imposes a significant annotation burden.

Consider a simple data transformation in Rust:

\begin{verbatim}
fn transform(data: Vec<i32>) -> Vec<i32> {
    data.into_iter().map(|x| x * 2).collect()
}
\end{verbatim}

The programmer must choose between \texttt{Vec<i32>} (ownership
transfer), \texttt{\&Vec<i32>} (immutable borrow), and
\texttt{\&mut Vec<i32>} (mutable borrow).  For a pure transformation,
this choice is an unnecessary cognitive burden.

In \JAPL{}, the same function requires no ownership annotations:

\begin{lstlisting}[caption={\JAPL{} vs.\ Rust: pure data transformation}]
fn transform(data: List[Int]) -> List[Int] =
  List.map(data, fn x -> x * 2)
\end{lstlisting}

The trade-off is that \JAPL{}'s pure layer requires a garbage
collector, while Rust avoids GC entirely.  We argue that this
trade-off is worthwhile for application-level code, where GC pauses
are negligible relative to I/O latency.

\subsection{Go: GC Everything}\label{sec:cmp-go}

Go takes the opposite extreme: all memory is garbage-collected, and
there is no ownership system.  Resources are managed by convention
(\texttt{defer} and the \texttt{io.Closer} interface) rather than by
the type system.

The consequence is that resource leaks are possible:

\begin{verbatim}
func processFile(path string) error {
    f, err := os.Open(path)
    if err != nil { return err }
    // Easy to forget: defer f.Close()
    data, err := io.ReadAll(f)
    // If we return here, f is leaked
    ...
}
\end{verbatim}

\JAPL{}'s linear types make this class of bug impossible: the
compiler rejects any program that fails to consume a resource.

\subsection{Haskell: The IO Monad}\label{sec:cmp-haskell}

Haskell confines side effects to the \texttt{IO} monad, ensuring that
pure functions are referentially transparent.  However, the monadic
style has several drawbacks:

\begin{enumerate}
  \item Monadic binding (\texttt{>>=}) is syntactically heavier than
        direct-style code.
  \item Composing different effects (state, error, IO) requires monad
        transformers, which are notoriously difficult to use.
  \item Resources within \texttt{IO} are not linearly tracked;
        finalizers and \texttt{bracket} provide safety but not
        compile-time guarantees.
\end{enumerate}

\JAPL{}'s effect system provides effect tracking without monadic
syntax: effects are listed after \texttt{with} in function signatures,
and composition is automatic (effects form a commutative monoid).
The resource layer adds linear tracking that Haskell's \texttt{IO}
lacks.

\subsection{Erlang: Copy Everything}\label{sec:cmp-erlang}

Erlang achieves concurrency safety by making all values immutable and
copying data between processes.  This eliminates shared mutable state
entirely but has performance implications for large data transfers.

\JAPL{} inherits Erlang's process isolation model for pure values but
extends it with ownership transfer for resources: a resource can be
moved between processes without copying, while maintaining the
invariant that only one process owns it at a time.

\subsection{OCaml: Mutable Refs}\label{sec:cmp-ocaml}

OCaml allows mutation through mutable record fields and
\texttt{ref} cells.  The type system does not distinguish pure and
impure code; any function may have side effects.

This pragmatic approach simplifies the type system but sacrifices
the ability to reason about purity.  \JAPL{}'s effect annotations
(\texttt{with Io}, \texttt{with State[s]}) restore this reasoning
while keeping the annotation burden low (effects are inferred within
function bodies).

\subsection{Clean: Uniqueness Types}\label{sec:cmp-clean}

Clean's uniqueness types~\cite{barendsen1996uniqueness} are the
closest relative of \JAPL{}'s resource layer.  Clean uses a
\texttt{*} attribute to mark unique values; the compiler ensures
that unique values are not aliased.

\JAPL{} extends Clean's approach in several ways:

\begin{enumerate}
  \item \textbf{Explicit layer separation}: \JAPL{} does not apply
        uniqueness to all types; only declared \texttt{resource} types
        are subject to linear discipline.
  \item \textbf{Borrowing}: Clean does not have a borrowing mechanism;
        in Clean, any use of a unique value consumes it.  \JAPL{}'s
        \texttt{ref} qualifier allows non-consuming reads.
  \item \textbf{Effect integration}: \JAPL{}'s resource operations
        are annotated with effects, providing an additional layer of
        safety.
  \item \textbf{Process integration}: \JAPL{} supports ownership
        transfer between processes, which Clean (not being an
        actor-based language) does not address.
\end{enumerate}

\subsection{Summary of Trade-offs}

\begin{table}[H]
\centering
\small
\caption{Trade-off analysis: safety guarantees vs.\ programmer burden}
\label{tab:tradeoffs}
\begin{tabular}{@{}lccccc@{}}
\toprule
& \rotatebox{60}{Data race free} & \rotatebox{60}{Resource safe}
& \rotatebox{60}{No GC needed} & \rotatebox{60}{Low annotation}
& \rotatebox{60}{Effect tracking} \\
\midrule
\JAPL{}   & \checkmark & \checkmark &            & \checkmark & \checkmark \\
Rust      & \checkmark & \checkmark & \checkmark &            &            \\
Go        &            &            &            & \checkmark &            \\
Haskell   & \checkmark &            &            &            & \checkmark \\
Erlang    & \checkmark &            &            & \checkmark &            \\
OCaml     &            &            &            & \checkmark &            \\
Clean     & \checkmark & \checkmark &            &            &            \\
\bottomrule
\end{tabular}
\end{table}

\JAPL{} is the only language in this comparison that simultaneously
provides data-race freedom, compile-time resource safety, low
annotation burden (for the majority of code), and effect tracking.
The cost is a garbage collector for the pure layer; we argue in
\Cref{sec:discussion} that this is an acceptable trade-off.


% ══════════════════════════════════════════════════════════════════
\section{Implementation}\label{sec:implementation}

\subsection{Architecture Overview}

The \JAPL{} compiler and runtime implement the dual-layer model as
follows:

\begin{enumerate}
  \item \textbf{Parser}: Produces an AST with explicit \texttt{use},
        \texttt{own}, and \texttt{ref} annotations.
  \item \textbf{Type checker}: Performs bidirectional type checking
        with local inference.  Maintains separate unrestricted
        ($\Gamma$) and linear ($\Delta$) contexts.
  \item \textbf{Linearity checker}: Verifies that all linear variables
        are consumed exactly once on every execution path.
  \item \textbf{Region inference}: Assigns regions to resource
        allocations and inserts release operations.
  \item \textbf{Uniqueness analysis}: Identifies pure values with a
        single reference for in-place update optimization.
  \item \textbf{Code generation}: Compiles to native code via LLVM or
        Cranelift.
  \item \textbf{Runtime}: Manages the dual-layer memory model at
        execution time.
\end{enumerate}

\subsection{Type Checker Implementation}\label{sec:impl-tc}

The type checker uses a \emph{splitting} judgment to partition the
linear context among subexpressions~\cite{walker2005substructural}:

\begin{definition}[Context splitting]\label{def:splitting}
A \emph{split} of linear context $\Delta$ into $\Delta_1$ and
$\Delta_2$ (written $\Delta = \Delta_1 \uplus \Delta_2$) assigns
each variable in $\Delta$ to exactly one of $\Delta_1$ or $\Delta_2$.
\end{definition}

For function application $e_1\; e_2$, the type checker:
\begin{enumerate}
  \item Generates a fresh split of the current linear context.
  \item Checks $e_1$ against $\Delta_1$.
  \item Checks $e_2$ against $\Delta_2$.
  \item Verifies that $\Delta_1$ and $\Delta_2$ are both fully consumed.
\end{enumerate}

For conditionals and pattern matching, the type checker verifies that
all branches consume the \emph{same} set of linear variables:

\begin{mathpar}
  \inferrule*[right=Match]
    {\Gamma; \Delta_0 \vdash e : \tau \\
     \forall i.\; \Gamma; \Delta_i \vdash p_i \Rightarrow e_i : \sigma \\
     \Delta_1 = \Delta_2 = \cdots = \Delta_n}
    {\Gamma; \Delta_0, \Delta_1 \vdash
      \textbf{match}\; e\; \textbf{with}\; \{p_i \Rightarrow e_i\} : \sigma}
\end{mathpar}

\subsection{Linearity Checking Algorithm}\label{sec:impl-linearity}

The linearity checker runs as a separate pass after type checking.
It traverses the typed AST and maintains a set of ``live'' linear
variables.  The algorithm is:

\begin{enumerate}
  \item At a \texttt{let} binding of a resource type: add the variable
        to the live set.
  \item At a use of a linear variable: remove it from the live set.
        If already removed, report a ``use after move'' error.
  \item At a branch point (if/match): check each branch independently;
        verify that all branches have the same live set at exit.
  \item At function exit: verify that the live set matches the function's
        return type (i.e., any linear variable still live must be
        returned).
  \item At scope exit for \texttt{use} bindings: if the variable is
        still live, insert a release call.
\end{enumerate}

\subsection{Runtime Memory Model}\label{sec:impl-runtime}

The runtime maintains two memory regions per process:

\paragraph{Immutable heap (GC-managed).}
All pure values reside here.  The garbage collector is generational
and per-process:

\begin{itemize}
  \item \textbf{Nursery}: Small, per-process bump allocator.  Most
        values die young and are collected cheaply.
  \item \textbf{Old generation}: Surviving values are promoted to a
        shared old generation, collected concurrently.
  \item \textbf{No write barriers}: Since data is immutable, there are
        no pointer mutations to track.  This significantly simplifies
        the GC.
  \item \textbf{Process death}: When a process terminates, its entire
        nursery is reclaimed instantly.
\end{itemize}

\paragraph{Resource arena (ownership-managed).}
All linear resources reside here.  There is no garbage collection;
resources are freed deterministically when consumed:

\begin{itemize}
  \item \textbf{Allocation}: Resources are allocated in a per-process
        arena.
  \item \textbf{Deallocation}: When a resource is consumed (released),
        its memory is immediately freed.
  \item \textbf{Transfer}: When ownership is transferred to another
        process, the resource is moved from one process's arena to
        another's.
  \item \textbf{No fragmentation}: Arena allocation avoids the
        fragmentation problems of general-purpose allocators.
\end{itemize}

\subsection{Compile-Time Ownership Verification}\label{sec:impl-verify}

The compiler performs ownership verification as a fixed-point
computation over the control-flow graph (CFG):

\begin{enumerate}
  \item Build the CFG for each function.
  \item At each CFG node, compute the set of live linear variables.
  \item At each join point, check that the live sets from all
        predecessors are identical.
  \item At the function exit node, check that the live set matches
        the expected output (linear variables that appear in the
        return type).
  \item If any check fails, emit a diagnostic with the offending
        variable and the point where it was consumed or escaped.
\end{enumerate}

This analysis is similar to liveness analysis in traditional compilers
but with the additional constraint that each linear variable must be
consumed \emph{exactly once}, not ``at least once'' or ``at most
once.''

\subsection{Performance Characteristics}

The dual-layer design has favorable performance characteristics:

\begin{itemize}
  \item \textbf{Pure-layer overhead}: GC pauses are bounded by the
        nursery size and are per-process (no global stop-the-world).
        Uniqueness inference eliminates unnecessary copies.
  \item \textbf{Resource-layer overhead}: Zero runtime overhead beyond
        the actual resource operations; ownership is verified entirely
        at compile time.
  \item \textbf{Combined}: The compiler can optimize across layers.
        For example, a buffer that is allocated, filled, and frozen
        within a single function can be stack-allocated if uniqueness
        analysis proves it does not escape.
\end{itemize}


% ══════════════════════════════════════════════════════════════════
\section{Formal Proofs}\label{sec:proofs}

This section proves type safety and resource safety for \JAPL{}'s
dual-layer system.  We work with a core calculus $\lambda^{\mathsf{JR}}$
(JAPL Resources) that captures the essential features.

\subsection{Syntax of $\lambda^{\mathsf{JR}}$}

\begin{definition}[Core calculus $\lambda^{\mathsf{JR}}$]\label{def:core-calculus}
\begin{align*}
  \text{Types}\quad \tau &::= B \mid \tau_1 \to \tau_2 \mid \tau_1 \times \tau_2
    \mid R \mid \Owned{R} \\
  \text{Expressions}\quad e &::= x \mid \lambda x.\, e \mid e_1\; e_2
    \mid (e_1, e_2) \mid \pi_i(e) \\
    &\phantom{::=} \mid \mathsf{new}_R \mid \mathsf{use}(e_1, e_2) \mid
      \mathsf{release}(e) \mid \mathsf{borrow}(e_1, e_2) \\
    &\phantom{::=} \mid \mathsf{freeze}(e) \mid v \\
  \text{Values}\quad v &::= \lambda x.\, e \mid (v_1, v_2) \mid c \mid r \\
  \text{Resources}\quad r &::= \mathsf{res}(a) \quad\text{(resource at address $a$)}
\end{align*}
where $B$ ranges over base types and $R$ ranges over resource type names.
\end{definition}

\begin{remark}[Effects in $\lambda^{\mathsf{JR}}$]\label{rem:effects-orthogonal}
The core calculus $\lambda^{\mathsf{JR}}$ deliberately omits effect
tracking ($\Io$, $\mathsf{State}[s]$, etc.) to isolate the
contribution of linear ownership.  Effect annotations appear in the
surface-language typing rules (\Cref{def:dual-judgment}) and in the
resource-specific rules (T-New, T-Release, T-Freeze), but a full
formalization of the effect system---including effect rows, handler
semantics, and the interaction between effect polymorphism and
linearity---is the subject of a companion paper on \JAPL{}'s
algebraic effect system.  The two systems compose orthogonally:
linearity governs \emph{how many times} a resource is used, while
effects govern \emph{what operations} may be performed.
\Cref{prop:effect-linear} states the key compatibility property
between the two.
\end{remark}

\begin{definition}[Store]\label{def:store}
A \emph{store} $\sigma$ is a partial function from addresses to resource
states:
\[
  \sigma : \mathsf{Addr} \rightharpoonup \mathsf{ResourceState}
\]
where $\mathsf{ResourceState} ::= \mathsf{Live}(v) \mid \mathsf{Released}$.
\end{definition}

\subsection{Operational Semantics}

We write $\bot_R$ for the \emph{default initial state} of resource
type $R$: the distinguished value that a freshly allocated resource
holds before any operations have been performed on it.  Each resource
type declaration implicitly defines its $\bot_R$ (e.g., an empty
buffer, an unconnected socket descriptor).

\begin{definition}[Small-step reduction]\label{def:reduction}
The reduction relation $\langle e, \sigma \rangle \longrightarrow
\langle e', \sigma' \rangle$ is defined by:

\begin{mathpar}
  \inferrule*[right=App]
    { }
    {\langle (\lambda x.\, e)\; v, \sigma \rangle \longrightarrow
     \langle [v/x]e, \sigma \rangle}

  \inferrule*[right=New]
    {a \notin \mathrm{dom}(\sigma)}
    {\langle \mathsf{new}_R, \sigma \rangle \longrightarrow
     \langle \mathsf{res}(a), \sigma[a \mapsto \mathsf{Live}(\bot_R)] \rangle}

  \inferrule*[right=Release]
    {\sigma(a) = \mathsf{Live}(v)}
    {\langle \mathsf{release}(\mathsf{res}(a)), \sigma \rangle \longrightarrow
     \langle (), \sigma[a \mapsto \mathsf{Released}] \rangle}

  \inferrule*[right=Freeze]
    {\sigma(a) = \mathsf{Live}(v)}
    {\langle \mathsf{freeze}(\mathsf{res}(a)), \sigma \rangle \longrightarrow
     \langle v, \sigma[a \mapsto \mathsf{Released}] \rangle}

  \inferrule*[right=Borrow]
    {\sigma(a) = \mathsf{Live}(v)}
    {\langle \mathsf{borrow}(\mathsf{res}(a), f), \sigma \rangle \longrightarrow
     \langle (f\; v, \mathsf{res}(a)), \sigma \rangle}
\end{mathpar}
\end{definition}

\subsection{Type Safety}

\begin{definition}[Store typing]\label{def:store-typing}
A store typing $\Sigma$ is a partial function from addresses to resource
types:
\[
  \Sigma : \mathsf{Addr} \rightharpoonup \{\mathsf{Live}(R) \mid
    \mathsf{Released}\}
\]
We write $\sigma : \Sigma$ when for every $a \in \mathrm{dom}(\Sigma)$:
\begin{itemize}
  \item If $\Sigma(a) = \mathsf{Live}(R)$, then $\sigma(a) = \mathsf{Live}(v)$
        for some $v$ of type $R$.
  \item If $\Sigma(a) = \mathsf{Released}$, then $\sigma(a) = \mathsf{Released}$.
\end{itemize}
\end{definition}

\begin{theorem}[Progress]\label{thm:progress}
If $\Gamma; \Delta \vdash_\varepsilon e : \tau$ and $\sigma : \Sigma$,
then either:
\begin{enumerate}[label=(\alph*)]
  \item $e$ is a value, or
  \item there exist $e'$, $\sigma'$ such that
        $\langle e, \sigma \rangle \longrightarrow
        \langle e', \sigma' \rangle$.
\end{enumerate}
\end{theorem}

\begin{proof}
By induction on the typing derivation $\Gamma; \Delta \vdash_\varepsilon
e : \tau$.

\textbf{Case} Var-Lin: $e = x$ and $x : \Owned{R} \in \Delta$.
Then $x$ is bound to a value (a resource reference $\mathsf{res}(a)$),
so $e$ is a value.

\textbf{Case} Var-Un: Similar.

\textbf{Case} $\lolli$-I: $e = \lambda x.\, e'$, which is a value.

\textbf{Case} $\lolli$-E: $e = e_1\; e_2$.  By induction, either
$e_1$ is a value or it can step.  If $e_1$ can step, so can $e$.
If $e_1$ is a value, it has type $A \lolli B$, so it is a lambda
$\lambda x.\, e'$.  By induction, either $e_2$ is a value or it can
step.  If $e_2$ can step, so can $e$.  If $e_2$ is a value $v$, then
$e$ can step by the App rule.

\textbf{Case} New: $e = \mathsf{new}_R$.  We can always find a fresh
address $a \notin \mathrm{dom}(\sigma)$ (addresses are unbounded), so
the New rule applies.

\textbf{Case} Release: $e = \mathsf{release}(\mathsf{res}(a))$.  By
the typing rule, $\mathsf{res}(a) : \Owned{R}$, so $a \in
\mathrm{dom}(\Sigma)$ with $\Sigma(a) = \mathsf{Live}(R)$.  Since
$\sigma : \Sigma$, we have $\sigma(a) = \mathsf{Live}(v)$, so the
Release rule applies.

The remaining cases (Freeze, Borrow, pairs, projections) follow
similarly.
\end{proof}

\begin{theorem}[Preservation]\label{thm:preservation}
If $\Gamma; \Delta \vdash_\varepsilon e : \tau$ and $\sigma : \Sigma$
and $\langle e, \sigma \rangle \longrightarrow \langle e', \sigma'
\rangle$, then there exist $\Gamma'$, $\Delta'$, $\Sigma'$ such that:
\begin{enumerate}[label=(\alph*)]
  \item $\Gamma'; \Delta' \vdash_\varepsilon e' : \tau$
  \item $\sigma' : \Sigma'$
  \item $\Sigma \subseteq \Sigma'$ (the store typing only grows or
        transitions $\mathsf{Live}$ entries to $\mathsf{Released}$)
\end{enumerate}
\end{theorem}

\begin{proof}
By induction on the typing derivation, with case analysis on the
reduction rule applied.

\textbf{Case} App: $e = (\lambda x.\, e')\; v$ steps to $[v/x]e'$.
By the T-App rule, the linear context is split as
$\Delta = \Delta_1 \uplus \Delta_2$ with
$\Gamma; \Delta_1 \vdash (\lambda x.\, e') : \tau_1 \lolli \tau$ and
$\Gamma; \Delta_2 \vdash v : \tau_1$.
By the substitution lemma for the dual-layer system (\Cref{lem:substitution}),
$\Gamma; \Delta' \vdash [v/x]e' : \tau$ where
$\Delta' = \Delta_1 \uplus (\Delta_2 \setminus \{x : \tau_1\})$,
i.e., $\Delta$ with $x$'s binding consumed.  The store is unchanged.

\textbf{Case} New: $e = \mathsf{new}_R$ steps to $\mathsf{res}(a)$
with $\sigma' = \sigma[a \mapsto \mathsf{Live}(\bot_R)]$.  Set
$\Sigma' = \Sigma[a \mapsto \mathsf{Live}(R)]$.  Then
$\mathsf{res}(a) : \Owned{R}$ under $\Delta' = \Delta, a : \Owned{R}$,
and $\sigma' : \Sigma'$.

\textbf{Case} Release: $e = \mathsf{release}(\mathsf{res}(a))$ steps
to $()$ with $\sigma' = \sigma[a \mapsto \mathsf{Released}]$.  Set
$\Sigma' = \Sigma[a \mapsto \mathsf{Released}]$.  The result $()$ has
type $\texttt{Unit}$.  The linear context $\Delta' = \Delta \setminus
\{a : \Owned{R}\}$: the binding is consumed.

\textbf{Case} Freeze: Similar to Release, but the result is the value
$v$ stored in the resource, which has a pure type.

\textbf{Case} Borrow: $e = \mathsf{borrow}(\mathsf{res}(a), f)$
steps to $(f\; v, \mathsf{res}(a))$.  The resource remains live; the
linear binding is preserved.  The result is a pair of the function
application result and the resource.

The remaining cases are standard.
\end{proof}

\begin{lemma}[Substitution]\label{lem:substitution}
If $\Gamma; \Delta, x : A \vdash_\varepsilon e : B$ and
$\Gamma; \cdot \vdash v : A$ (if $A$ is unrestricted) or
$v = \mathsf{res}(a)$ with $a : \Owned{R}$ (if $A$ is linear),
then $\Gamma; \Delta' \vdash_\varepsilon [v/x]e : B$ where:
\begin{itemize}
  \item $\Delta' = \Delta$ if $A$ is unrestricted.
  \item $\Delta' = \Delta$ if $A$ is linear and $x$ is consumed exactly
        once in $e$.
\end{itemize}
\end{lemma}

\begin{proof}
By induction on the typing derivation, tracking the number of
occurrences of $x$ in $e$.  The linearity checker ensures exactly one
occurrence for linear variables.
\end{proof}

\subsection{Resource Safety}

\begin{theorem}[Resource safety]\label{thm:resource-safety-proof}
If $\Gamma; \cdot \vdash_\varepsilon e : \tau$ (the program type-checks
with an empty linear context) and
$\langle e, \emptyset \rangle \longrightarrow^* \langle v, \sigma \rangle$
(the program terminates with value $v$ and final store $\sigma$), then
for every $a \in \mathrm{dom}(\sigma)$:
\[
  \sigma(a) = \mathsf{Released}
\]
That is, every resource acquired during execution has been released.
\end{theorem}

\begin{proof}
We establish the invariant: at every step of reduction, the number of
$\mathsf{Live}$ entries in $\sigma$ equals the size of the linear
context $\Delta$.

\textbf{Base case}: Initially, $\sigma = \emptyset$ and $\Delta =
\cdot$, so the invariant holds trivially.

\textbf{Inductive step}: By preservation (\Cref{thm:preservation}),
each reduction step preserves the correspondence between $\Delta$ and
$\Sigma$.  Specifically:
\begin{itemize}
  \item New: adds one entry to both $\sigma$ ($\mathsf{Live}$) and
        $\Delta$, preserving the count.
  \item Release/Freeze: transitions one entry in $\sigma$ from
        $\mathsf{Live}$ to $\mathsf{Released}$ and removes one entry
        from $\Delta$, preserving the count.
  \item All other rules: preserve both $\sigma$ and the size of $\Delta$.
\end{itemize}

\textbf{Conclusion}: When the program terminates with
$\Gamma; \cdot \vdash v : \tau$ (empty linear context), the number
of $\mathsf{Live}$ entries in $\sigma$ is zero.  Therefore every
entry is $\mathsf{Released}$.
\end{proof}

\begin{corollary}[No resource leaks]\label{cor:no-leaks}
Well-typed \JAPL{} programs (without \texttt{unsafe}) cannot leak
resources: every file is closed, every socket is disconnected, every
buffer is freed.
\end{corollary}

\begin{corollary}[No use-after-free]\label{cor:no-uaf}
Well-typed \JAPL{} programs cannot use a resource after it has been
released: the linear type system ensures that the binding is consumed
by the release operation and cannot be used again.
\end{corollary}

\begin{corollary}[No double-free]\label{cor:no-double-free}
Well-typed \JAPL{} programs cannot release a resource twice: the linear
type system ensures exactly-once consumption.
\end{corollary}


% ══════════════════════════════════════════════════════════════════
\section{Discussion}\label{sec:discussion}

\subsection{When Ownership Pressure Leaks}

The dual-layer design succeeds when the boundary between pure and
resource code is clean.  In practice, however, there are cases where
ownership concerns propagate into otherwise-pure code:

\paragraph{Resource-carrying data structures.}
If a data structure contains a resource (e.g., a list of open file
handles), the entire structure becomes linear.  This forces the
programmer to manage ownership for the container, even if the
business logic only cares about the data inside.

\paragraph{Callbacks with resources.}
A higher-order function that accepts a callback may need to thread
resources through the callback, complicating the interface.

\paragraph{Long-lived resources.}
A database connection pool manages resources with lifetimes that span
many function calls.  The pool itself becomes a resource, and every
function that uses the pool must participate in the ownership protocol.

\subsection{Mitigations}

\JAPL{} provides several mechanisms to reduce ownership pressure:

\paragraph{Borrowing (\texttt{ref}).}
Most pure functions that need to read a resource can take a borrowed
reference, avoiding ownership transfer entirely.

\paragraph{Scoped resources (\texttt{use}).}
The \texttt{use} binding automatically manages resource lifetimes
within a scope, reducing boilerplate.

\paragraph{Process encapsulation.}
Long-lived resources can be owned by a dedicated process, which
exposes a pure message-passing interface.  Other processes interact
with the resource through messages, never seeing the ownership:

\begin{lstlisting}[caption={Process encapsulation hides ownership},label={lst:encapsulation}]
-- The database process owns the connection
fn db_process(conn: Owned[DbConnection])
    -> Never with Io, Process[DbMsg] =
  match Process.receive() with
  | Query(sql, reply) ->
      let (conn, result) = Db.query(conn, sql)
      Reply.send(reply, result)
      db_process(conn)
  | Shutdown ->
      Db.close(conn)
      Process.exit(Normal)

-- Callers never see Owned[DbConnection]
fn get_user(db: Pid[DbMsg], id: UserId)
    -> Result[User, DbError] with Process =
  let result = Process.call(db, fn reply -> Query(user_sql(id), reply))
  parse_user_row(result)
\end{lstlisting}

This pattern is the recommended way to manage long-lived resources
in \JAPL{}: the process boundary naturally encapsulates the
ownership, and the caller's interface is pure message-passing.

\paragraph{Region-based grouping.}
Multiple resources acquired in the same scope can be grouped into a
single region, reducing the number of individual ownership threads.

\subsection{The 80/20 of Resource Management}

Our experience with prototype \JAPL{} applications suggests the
following distribution:

\begin{itemize}
  \item \textbf{80\%} of code is pure: data transformations, business
        logic, protocol handling, serialization.  No ownership
        annotations needed.
  \item \textbf{15\%} of code uses scoped resources (\texttt{use}):
        open file, read, close; connect, query, disconnect.  Ownership
        is automatic via the \texttt{use} binding.
  \item \textbf{5\%} of code requires explicit ownership management:
        long-lived resources, cross-process transfers, complex
        resource graphs.  This is where the linear type system
        provides its greatest value.
\end{itemize}

The dual-layer design ensures that the 80\% majority pays no
annotation cost, while the 5\% minority gets strong compile-time
guarantees.

\subsection{Effect System Synergy}

The effect system and the ownership system reinforce each other.  The
effect system ensures that resource operations are visible in function
signatures (\texttt{with Io}), preventing hidden side effects.  The
ownership system ensures that resources are properly managed even when
effects are composed.

Together, they provide a stronger guarantee than either alone: a
function with signature \texttt{fn(x: Int) -> Int} is guaranteed to
be pure---it cannot secretly open a file, allocate a buffer, or
access the network.

\subsection{Comparison with Monadic IO}

The Haskell approach confines \emph{all} effects to the \texttt{IO}
monad, including both resource management and general side effects.
This conflates two concerns:

\begin{enumerate}
  \item \textbf{Effect tracking}: Knowing that a function performs IO.
  \item \textbf{Resource management}: Ensuring that resources are
        properly released.
\end{enumerate}

\JAPL{} separates these concerns: the effect system handles (1), and
the ownership system handles (2).  This separation allows finer-grained
reasoning: two functions may both have effect \texttt{Io}, but one
manages resources (visible in its type via \texttt{Owned}) while the
other merely reads a file.

\subsection{Limitations}

The dual-layer approach has several limitations:

\begin{enumerate}
  \item \textbf{GC dependency}: The pure layer requires a garbage
        collector, making \JAPL{} unsuitable for environments where
        GC is unacceptable (embedded systems, real-time audio).
  \item \textbf{Layer crossing cost}: Moving data between layers
        (freezing a buffer into immutable bytes) may involve a copy
        if the resource is not uniquely held.
  \item \textbf{Learning curve}: Programmers must understand the
        distinction between pure and resource types, even if the
        pure layer requires no annotations.
  \item \textbf{Expressiveness}: Some patterns that are natural in
        Rust (e.g., self-referential structures with lifetimes) are
        more difficult to express in \JAPL{}'s simpler ownership model.
\end{enumerate}

\subsection{Future Work}

Several directions remain for future investigation:

\begin{enumerate}
  \item \textbf{Graded linearity}: Allowing resources to be used
        exactly $n$ times (useful for connection pools with bounded
        concurrency).
  \item \textbf{Session types}: Extending resource types with session
        type protocols, ensuring that resources are used in the
        correct order (e.g., connect, authenticate, query, disconnect).
  \item \textbf{Dependent resource types}: Allowing resource types to
        depend on values, enabling richer invariants (e.g., a file
        handle that tracks whether it is open for reading, writing,
        or both).
  \item \textbf{Formal verification}: Connecting \JAPL{}'s type system
        to separation logic via the LNL model, enabling machine-checked
        proofs of resource safety.
\end{enumerate}


% ══════════════════════════════════════════════════════════════════
\section{Conclusion}\label{sec:conclusion}

We have presented \JAPL{}'s second core principle: \textbf{mutation is
local and explicit}.  The dual-layer type system---pure values managed
by GC, linear resources managed by ownership---achieves a design point
that has not been explored by prior languages:

\begin{enumerate}
  \item \textbf{Safety without pervasive annotation}: The majority of
        code lives in the pure layer and requires no ownership
        annotations, lifetime parameters, or borrow-checker reasoning.
  \item \textbf{Resource safety by construction}: Every resource
        acquired by a well-typed program is released exactly once,
        as proven by our type safety and resource safety theorems.
  \item \textbf{Effect visibility}: The effect system makes mutation
        visible in function signatures, enabling equational reasoning
        for pure code and local reasoning for effectful code.
  \item \textbf{Process integration}: Ownership transfer between
        processes extends the linear discipline to concurrent settings,
        while process encapsulation hides ownership behind pure
        message-passing interfaces.
\end{enumerate}

The formal framework---rooted in Girard's linear logic, Benton's
LNL adjunction model, and the theory of substructural type
systems---provides a rigorous foundation for the design.  The
categorical interpretation (\Cref{sec:categorical}) shows that
\JAPL{}'s dual-layer system is not an \emph{ad hoc} engineering
compromise but a principled instantiation of the well-studied
relationship between linear and Cartesian closed categories.

We believe that the dual-layer approach represents a practical sweet
spot for language design: it gives systems programmers the control they
need over resources, while giving application programmers the
simplicity they want for everyday code.  The principle that mutation
should be local and explicit---not absent, not pervasive, but
\emph{controlled}---is central to \JAPL{}'s mission of being ``pure
by default, concurrent by design, resource-safe by construction,
distributed without apology.''


% ══════════════════════════════════════════════════════════════════
\bibliographystyle{plainnat}

\begin{thebibliography}{50}

\bibitem[Girard(1987)]{girard1987linear}
Jean-Yves Girard.
\newblock Linear logic.
\newblock \emph{Theoretical Computer Science}, 50(1):1--102, 1987.

\bibitem[Wadler(1990)]{wadler1990linear}
Philip Wadler.
\newblock Linear types can change the world!
\newblock In \emph{Programming Concepts and Methods}, pages 561--581. North-Holland, 1990.

\bibitem[Wadler(1993)]{wadler1993taste}
Philip Wadler.
\newblock A taste of linear logic.
\newblock In \emph{Mathematical Foundations of Computer Science}, volume 711 of LNCS, pages 185--210. Springer, 1993.

\bibitem[Wadler(2012)]{wadler2012propositions}
Philip Wadler.
\newblock Propositions as sessions.
\newblock In \emph{Proceedings of the 17th ACM SIGPLAN International Conference on Functional Programming (ICFP)}, pages 273--286, 2012.

\bibitem[Walker(2005)]{walker2005substructural}
David Walker.
\newblock Substructural type systems.
\newblock In Benjamin~C. Pierce, editor, \emph{Advanced Topics in Types and Programming Languages}, chapter~1, pages 3--43. MIT Press, 2005.

\bibitem[Pierce(2002)]{pierce2002types}
Benjamin~C. Pierce.
\newblock \emph{Types and Programming Languages}.
\newblock MIT Press, 2002.

\bibitem[Pierce(2004)]{pierce2004advanced}
Benjamin~C. Pierce, editor.
\newblock \emph{Advanced Topics in Types and Programming Languages}.
\newblock MIT Press, 2004.

\bibitem[Benton(1995)]{benton1995mixed}
P.~Nick Benton.
\newblock A mixed linear and non-linear logic: Proofs, terms and models.
\newblock In \emph{Computer Science Logic (CSL)}, volume 933 of LNCS, pages 121--135. Springer, 1995.

\bibitem[Barr(1991)]{barr1991autonomous}
Michael Barr.
\newblock {$\ast$}-Autonomous categories and linear logic.
\newblock \emph{Mathematical Structures in Computer Science}, 1(2):159--178, 1991.

\bibitem[Bernardy et~al.(2018)]{bernardy2018linear}
Jean-Philippe Bernardy, Mathieu Boespflug, Ryan~R. Newton, Simon Peyton~Jones, and Arnaud Spiwack.
\newblock Linear {H}askell: Practical linearity in a higher-order polymorphic language.
\newblock \emph{Proceedings of the ACM on Programming Languages}, 2(POPL):1--29, 2018.

\bibitem[Rust~Team(2021)]{rust2021}
The Rust~Team.
\newblock \emph{The Rust Programming Language}.
\newblock No Starch Press, 2021.
\newblock \url{https://doc.rust-lang.org/book/}.

\bibitem[Jung et~al.(2017)]{jung2017rustbelt}
Ralf Jung, Jacques-Henri Jourdan, Robbert Krebbers, and Derek Dreyer.
\newblock {RustBelt}: Securing the foundations of the {R}ust programming language.
\newblock \emph{Proceedings of the ACM on Programming Languages}, 2(POPL):1--34, 2017.

\bibitem[Barendsen and Smetsers(1996)]{barendsen1996uniqueness}
Erik Barendsen and Sjaak Smetsers.
\newblock Uniqueness typing for functional languages with graph rewriting semantics.
\newblock \emph{Mathematical Structures in Computer Science}, 6(6):579--612, 1996.

\bibitem[Plasmeijer and van Eekelen(2001)]{plasmeijer2001clean}
Rinus Plasmeijer and Marko van Eekelen.
\newblock \emph{Clean Language Report, version 2.1}.
\newblock University of Nijmegen, 2001.

\bibitem[Tofte and Talpin(1997)]{tofte1997region}
Mads Tofte and Jean-Pierre Talpin.
\newblock Region-based memory management.
\newblock \emph{Information and Computation}, 132(2):109--176, 1997.

\bibitem[Jim et~al.(2002)]{jim2002cyclone}
Trevor Jim, J.~Gregory Morrisett, Dan Grossman, Michael~W. Hicks, James Cheney, and Yanling Wang.
\newblock {C}yclone: A safe dialect of {C}.
\newblock In \emph{USENIX Annual Technical Conference}, pages 275--288, 2002.

\bibitem[Grossman et~al.(2002)]{grossman2002region}
Dan Grossman, Greg Morrisett, Trevor Jim, Michael Hicks, Yanling Wang, and James Cheney.
\newblock Region-based memory management in {C}yclone.
\newblock In \emph{Proceedings of the ACM SIGPLAN Conference on Programming Language Design and Implementation (PLDI)}, pages 282--293, 2002.

\bibitem[Dennis and Van~Horn(1966)]{dennis1966programming}
Jack~B. Dennis and Earl~C. Van~Horn.
\newblock Programming semantics for multiprogrammed computations.
\newblock \emph{Communications of the ACM}, 9(3):143--155, 1966.

\bibitem[Miller(2006)]{miller2006robust}
Mark~S. Miller.
\newblock \emph{Robust Composition: Towards a Unified Approach to Access Control and Concurrency Control}.
\newblock PhD thesis, Johns Hopkins University, 2006.

\bibitem[Plotkin and Pretnar(2009)]{plotkin2009handlers}
Gordon Plotkin and Matija Pretnar.
\newblock Handlers of algebraic effects.
\newblock In \emph{Programming Languages and Systems (ESOP)}, volume 5502 of LNCS, pages 80--94. Springer, 2009.

\bibitem[Bauer and Pretnar(2015)]{bauer2015programming}
Andrej Bauer and Matija Pretnar.
\newblock Programming with algebraic effects and handlers.
\newblock \emph{Journal of Logical and Algebraic Methods in Programming}, 84(1):108--123, 2015.

\bibitem[Kammar et~al.(2013)]{kammar2013handlers}
Ohad Kammar, Sam Lindley, and Nicolas Oury.
\newblock Handlers in action.
\newblock In \emph{Proceedings of the 18th ACM SIGPLAN International Conference on Functional Programming (ICFP)}, pages 145--158, 2013.

\bibitem[Wright and Felleisen(1994)]{wright1994syntactic}
Andrew~K. Wright and Matthias Felleisen.
\newblock A syntactic approach to type soundness.
\newblock \emph{Information and Computation}, 115(1):38--94, 1994.

\bibitem[Armstrong(2003)]{armstrong2003erlang}
Joe Armstrong.
\newblock \emph{Programming Erlang: Software for a Concurrent World}.
\newblock Pragmatic Bookshelf, 2003.

\bibitem[Peyton~Jones(2003)]{peytonjones2003haskell}
Simon Peyton~Jones.
\newblock \emph{Haskell 98 Language and Libraries: The Revised Report}.
\newblock Cambridge University Press, 2003.

\bibitem[Leroy(2000)]{leroy2000ocaml}
Xavier Leroy.
\newblock \emph{The {OCaml} System: Documentation and User's Manual}.
\newblock INRIA, 2000.

\bibitem[Donovan and Kernighan(2015)]{donovan2015go}
Alan A.~A. Donovan and Brian~W. Kernighan.
\newblock \emph{The Go Programming Language}.
\newblock Addison-Wesley, 2015.

\bibitem[Clebsch et~al.(2015)]{clebsch2015deny}
Sylvan Clebsch, Sophia Drossopoulou, Sebastian Blessing, and Andy McNeil.
\newblock Deny capabilities for safe, fast actors.
\newblock In \emph{Proceedings of the 5th International Workshop on Programming Based on Actors, Agents, and Decentralized Control}, pages 1--12, 2015.

\bibitem[Tov and Pucella(2011)]{tov2011practical}
Jesse~A. Tov and Riccardo Pucella.
\newblock Practical affine types.
\newblock In \emph{Proceedings of the 38th ACM SIGPLAN-SIGACT Symposium on Principles of Programming Languages (POPL)}, pages 447--458, 2011.

\bibitem[Leijen(2017)]{leijen2017koka}
Daan Leijen.
\newblock Type directed compilation of row-typed algebraic effects.
\newblock In \emph{Proceedings of the 44th ACM SIGPLAN Symposium on Principles of Programming Languages (POPL)}, pages 486--499, 2017.

\bibitem[Mazurak et~al.(2010)]{mazurak2010lightweight}
Karl Mazurak, Jianzhou Zhao, and Steve Zdancewic.
\newblock Lightweight linear types in {S}ystem {F}$^\circ$.
\newblock In \emph{Proceedings of the 5th ACM SIGPLAN Workshop on Types in Language Design and Implementation (TLDI)}, pages 77--88, 2010.

\bibitem[Krishnaswami et~al.(2019)]{krishnaswami2019integrating}
Neel Krishnaswami, Pierre Pradic, and Nick Benton.
\newblock Integrating linear and dependent types.
\newblock In \emph{Proceedings of the ACM on Programming Languages}, 3(POPL):1--28, 2019.

\bibitem[Sewell et~al.(2010)]{sewell2010ott}
Peter Sewell, Francesco~Zappa Nardelli, Scott Owens, Gilles Peskine, Thomas Ridge, Susmit Sarkar, and Rok Strni\v{s}a.
\newblock Ott: Effective tool support for the working semanticist.
\newblock \emph{Journal of Functional Programming}, 20(1):71--122, 2010.

\bibitem[Milner(1978)]{milner1978theory}
Robin Milner.
\newblock A theory of type polymorphism in programming.
\newblock \emph{Journal of Computer and System Sciences}, 17(3):348--375, 1978.

\bibitem[Reynolds(1983)]{reynolds1983types}
John~C. Reynolds.
\newblock Types, abstraction and parametric polymorphism.
\newblock In \emph{Information Processing 83}, pages 513--523. Elsevier, 1983.

\bibitem[Moggi(1991)]{moggi1991notions}
Eugenio Moggi.
\newblock Notions of computation and monads.
\newblock \emph{Information and Computation}, 93(1):55--92, 1991.

\bibitem[Austral(2023)]{austral2023}
Fernando Borretti.
\newblock The {A}ustral programming language.
\newblock \url{https://austral-lang.org/}, 2023.

\bibitem[Abramsky(1993)]{abramsky1993computational}
Samson Abramsky.
\newblock Computational interpretations of linear logic.
\newblock \emph{Theoretical Computer Science}, 111(1--2):3--57, 1993.

\bibitem[Gaboardi et~al.(2016)]{gaboardi2016combining}
Marco Gaboardi, Shin-ya Katsumata, Dominic Orchard, Flavien Breuvart, and Tarmo Uustalu.
\newblock Combining effects and coeffects via grading.
\newblock In \emph{Proceedings of the 21st ACM SIGPLAN International Conference on Functional Programming (ICFP)}, pages 476--489, 2016.

\end{thebibliography}

% ══════════════════════════════════════════════════════════════════
\appendix

\section{Complete Typing Rules for $\lambda^{\mathsf{JR}}$}\label{app:rules}

For reference, we present the complete set of typing rules for the
core calculus.

\begin{definition}[Full typing rules]\label{def:full-rules}
~

\begin{mathpar}
  \inferrule*[right=T-Var-Un]
    {x : \tau \in \Gamma}
    {\Gamma; \cdot \vdash x : \tau}

  \inferrule*[right=T-Var-Lin]
    { }
    {\Gamma; x : \tau \vdash x : \tau}

  \inferrule*[right=T-Abs]
    {\Gamma; \Delta, x : \tau_1 \vdash e : \tau_2}
    {\Gamma; \Delta \vdash \lambda x.\, e : \tau_1 \lolli \tau_2}

  \inferrule*[right=T-Abs-Un]
    {\Gamma, x : \tau_1; \cdot \vdash e : \tau_2}
    {\Gamma; \cdot \vdash \lambda x.\, e : \tau_1 \to \tau_2}

  \inferrule*[right=T-App]
    {\Gamma; \Delta_1 \vdash e_1 : \tau_1 \lolli \tau_2 \\
     \Gamma; \Delta_2 \vdash e_2 : \tau_1}
    {\Gamma; \Delta_1, \Delta_2 \vdash e_1\; e_2 : \tau_2}

  \inferrule*[right=T-App-Un]
    {\Gamma; \Delta_1 \vdash e_1 : \tau_1 \to \tau_2 \\
     \Gamma; \Delta_2 \vdash e_2 : \tau_1}
    {\Gamma; \Delta_1, \Delta_2 \vdash e_1\; e_2 : \tau_2}

  \inferrule*[right=T-Pair]
    {\Gamma; \Delta_1 \vdash e_1 : \tau_1 \\
     \Gamma; \Delta_2 \vdash e_2 : \tau_2}
    {\Gamma; \Delta_1, \Delta_2 \vdash (e_1, e_2) : \tau_1 \tensor \tau_2}

  \inferrule*[right=T-LetPair]
    {\Gamma; \Delta_1 \vdash e_1 : \tau_1 \tensor \tau_2 \\
     \Gamma; \Delta_2, x : \tau_1, y : \tau_2 \vdash e_2 : \tau}
    {\Gamma; \Delta_1, \Delta_2 \vdash
      \mathbf{let}\; (x, y) = e_1\; \mathbf{in}\; e_2 : \tau}

  \inferrule*[right=T-New]
    { }
    {\Gamma; \cdot \vdash_{\Io} \mathsf{new}_R : \Owned{R}}

  \inferrule*[right=T-Release]
    {\Gamma; \Delta \vdash e : \Owned{R}}
    {\Gamma; \Delta \vdash_{\Io} \mathsf{release}(e) : \mathsf{Unit}}

  \inferrule*[right=T-Freeze]
    {\Gamma; \Delta \vdash e : \Owned{R} \\
     R \hookrightarrow \tau}
    {\Gamma; \Delta \vdash_{\Io} \mathsf{freeze}(e) : \tau}

  \inferrule*[right=T-Borrow]
    {\Gamma; \Delta, x : \Owned{R} \vdash e_1 : \Owned{R} \\
     \Gamma; y : \rref\; R \vdash e_2 : \tau}
    {\Gamma; \Delta, x : \Owned{R} \vdash
      \mathsf{borrow}(x\; \mathbf{as}\; y,\; e_2) : \tau \tensor \Owned{R}}

  \inferrule*[right=T-If]
    {\Gamma; \Delta_0 \vdash e : \mathsf{Bool} \\
     \Gamma; \Delta \vdash e_1 : \tau \\
     \Gamma; \Delta \vdash e_2 : \tau}
    {\Gamma; \Delta_0, \Delta \vdash
      \mathbf{if}\; e\; \mathbf{then}\; e_1\; \mathbf{else}\; e_2 : \tau}
\end{mathpar}

Note that T-Abs-Un requires the linear context to be empty ($\cdot$):
an unrestricted function $\tau_1 \to \tau_2$ may be duplicated, so it
must not close over linear resources---otherwise, applying the
closure twice would duplicate those resources, violating linearity.

Note that T-If requires both branches to consume the \emph{same}
linear context $\Delta$.  This is the key rule that prevents
conditional resource leaks.
\end{definition}

\section{Extended Examples}\label{app:examples}

\subsection{Database Connection Pool}

\begin{lstlisting}[caption={A resource-safe connection pool},label={lst:pool}]
resource DbConnection

type PoolMsg =
  | Checkout(Reply[Owned[DbConnection]])
  | Return(Owned[DbConnection])
  | Shutdown

fn pool_process(conns: List[Owned[DbConnection]],
                waiting: List[Reply[Owned[DbConnection]]])
    -> Never with Io, Process[PoolMsg] =
  match Process.receive() with
  | Checkout(reply) ->
      match conns with
      | [conn, ..rest] ->
          Reply.send(reply, conn)
          pool_process(rest, waiting)
      | [] ->
          pool_process(conns, List.append(waiting, [reply]))
  | Return(conn) ->
      match waiting with
      | [waiter, ..rest] ->
          Reply.send(waiter, conn)
          pool_process(conns, rest)
      | [] ->
          pool_process([conn, ..conns], waiting)
  | Shutdown ->
      List.each(conns, fn conn -> Db.close(conn))
      Process.exit(Normal)
\end{lstlisting}

\subsection{State Effect with Linear Resources}

\begin{lstlisting}[caption={Composing effects with linear resources},label={lst:effect-compose}]
-- Explicit mutable state via effect
fn accumulate(items: List[Int]) -> Int with State[Int] =
  List.each(items, fn x -> State.modify(fn acc -> acc + x))
  State.get()

-- Resource layer: mutable buffer (ownership tracked)
fn fill_buffer(data: Bytes) -> Bytes with Io =
  let buf = Buffer.alloc(1024)
  let buf = Buffer.write(buf, 0, data)
  Buffer.freeze(buf)

-- Combining pure, state, and resource layers
fn process_and_store(items: List[Int], path: Path)
    -> Result[Unit, IoError] with Io =
  let total = State.run(0, fn -> accumulate(items))
  let file = File.open(path, Write)?
  File.write(file, Int.to_string(total))
  File.close(file)
  Ok(())
\end{lstlisting}

\subsection{Safe FFI Wrapping}

\begin{lstlisting}[caption={Safe wrapper around a C FFI resource},label={lst:ffi-wrap}]
-- Foreign resource: a C library handle
resource CHandle

foreign "C" fn c_lib_open(config: CString) -> Ptr[CHandle]
foreign "C" fn c_lib_process(h: Ptr[CHandle], data: Ptr[Byte], len: CInt) -> CInt
foreign "C" fn c_lib_close(h: Ptr[CHandle]) -> Unit

-- Safe JAPL wrapper: all ownership is explicit
fn open_handle(config: String) -> Result[Owned[CHandle], Error] with Io =
  use cstr = CString.from(config)
  let ptr = unsafe c_lib_open(cstr)
  if Ptr.is_null(ptr) then Err(OpenFailed)
  else Ok(CHandle.from_raw(ptr))

fn process_data(handle: Owned[CHandle], data: Bytes)
    -> (Owned[CHandle], Result[Int, Error]) with Io =
  let result = unsafe c_lib_process(
    CHandle.as_ptr(ref handle),
    Bytes.as_ptr(data),
    Bytes.length(data)
  )
  if result < 0 then (handle, Err(ProcessFailed(result)))
  else (handle, Ok(result))

fn close_handle(handle: Owned[CHandle]) -> Unit with Io =
  unsafe c_lib_close(CHandle.into_raw(handle))
\end{lstlisting}

\end{document}
